%%%%%%%%%%%%%%%%%%%%%%%%%%%%%%%%%%%%%%%%%%%%%%%%%%%%%%%%%%%%%%%%%%%%%%%%%%%%%%%
%% sec_theory_current.tex
%%
%% Section 2.5 for paper2_main.tex
%% Title: Current Operator and Nonlocal Pseudopotential Contributions
%%
%% %%%%%%%%%%%%%%%%%%%%%%%%%%%%%%%%%%%%%%%%%%%%%%%%%%%%%%%%%%%%%%%%%%%%%%%%%%%%%%%
%% sec_theory_current.tex
%%
%% Section 2.5 for paper2_main.tex
%% Title: Current Operator and Nonlocal Pseudopotential Contributions
%%
%% %%%%%%%%%%%%%%%%%%%%%%%%%%%%%%%%%%%%%%%%%%%%%%%%%%%%%%%%%%%%%%%%%%%%%%%%%%%%%%%
%% sec_theory_current.tex
%%
%% Section 2.5 for paper2_main.tex
%% Title: Current Operator and Nonlocal Pseudopotential Contributions
%%
%% %%%%%%%%%%%%%%%%%%%%%%%%%%%%%%%%%%%%%%%%%%%%%%%%%%%%%%%%%%%%%%%%%%%%%%%%%%%%%%%
%% sec_theory_current.tex
%%
%% Section 2.5 for paper2_main.tex
%% Title: Current Operator and Nonlocal Pseudopotential Contributions
%%
%% \input{sec_theory_current} replaces the placeholder in sec_theory.tex.
%%
%% CONVENTIONS:
%%   Atomic units: hbar = m_e = e = 4*pi*eps_0 = 1
%%   Electron charge = -1  (e > 0 is the elementary charge magnitude)
%%   In atomic units the Hamiltonian is
%%       H = (1/2)(-i*nabla + A)^2 + V_loc + V_nl
%%   Imaginary unit: \imath  (dotless i, consistent with paper 1)
%%   Sign check: NEAT-notes-RMG has a typo (-1/2m kinetic energy);
%%               NEAT-project-notes has the correct +1/2m sign — used here.
%%
%% NOTATION MACROS (all in notation.tex):
%%   \phat   = \hat{\mathbf{p}}  canonical momentum
%%   \pihat  = \hat{\boldsymbol{\pi}}  mechanical momentum
%%   \vhat   = \hat{\mathbf{v}}  velocity operator
%%   \rhat   = \hat{\mathbf{r}}  position
%%   \Vnl    = V^{\mathrm{nl}}
%%   \Vloc   = V^{\mathrm{loc}}
%%   \Avec   = \mathbf{A}
%%   \Jpara  = \mathbf{J}_{\mathrm{para}}
%%   \Jdia   = \mathbf{J}_{\mathrm{dia}}
%%   \Jmat{x}= \mathcal{J}^{x}
%%%%%%%%%%%%%%%%%%%%%%%%%%%%%%%%%%%%%%%%%%%%%%%%%%%%%%%%%%%%%%%%%%%%%%%%%%%%%%%

\subsection{Current Operator and Nonlocal Pseudopotential Contributions}
\label{sec:current}

The central observable in periodic RT-TDDFT is the macroscopic
electronic current density $J_\alpha(t)$, from which the optical
conductivity $\sigma_{\alpha\beta}(\omega)$ and dielectric function
$\varepsilon_{\alpha\beta}(\omega)$ are obtained by Fourier analysis
(Section~\ref{sec:postprocessing}).
For periodic systems the electric dipole moment
$\langle\hat{\mathbf{r}}\rangle$ is ill-defined because the position
operator is incompatible with periodic boundary
conditions,\cite{resta-1998,king-smith-vanderbilt-1993} and the
current is the natural replacement.

We derive the current density operator in four steps: (i) we
distinguish canonical from mechanical momentum and define the velocity
operator via the Heisenberg equation of motion, identifying the role
of nonlocal pseudopotentials; (ii) we rederive the same result from
the continuity equation, decomposing the current into paramagnetic,
diamagnetic, and nonlocal pseudopotential contributions; (iii) we
specialize to Bloch states and derive the $\mathbf{k}$-vector
correction explicitly from the Bloch ansatz; (iv) we express the
macroscopic current in density matrix form and verify the sign
convention in the code.
All equations are in atomic units ($\hbar = m_e = e = 4\pi\varepsilon_0
= 1$) with the electron charge equal to $-1$.

%--------------------------------------------------------------------
\subsubsection{Canonical and Mechanical Momentum}
\label{sec:canon-mech}
%--------------------------------------------------------------------

The Hamiltonian of the Kohn-Sham system in the presence of a
spatially uniform, time-dependent vector potential $\Avec(t)$
(long-wavelength approximation, see Section~\ref{sec:EM-coupling})
is
\begin{equation}
  \hat{H}(t) = \frac{1}{2}
    \!\left(\hat{\mathbf{p}} + \Avec(t)\right)^{\!2}
    + \Vloc(\mathbf{r}) + \Vnl,
  \label{eq:H-vg-2}
\end{equation}
where $\hat{\mathbf{p}} = -\imath\nabla$ is the \emph{canonical}
momentum operator and the sign on $\Avec$ reflects that in atomic
units with electron charge $-1$ the coupling is
$q\Avec = (+1)\cdot(-1)\cdot\Avec = -\Avec$, giving mechanical
momentum $\pihat = \hat{\mathbf{p}} + \Avec$.\footnote{In
Gaussian units the coupling is $q\Avec/c = -e\Avec/c$ for an
electron; in atomic units $c=1/\alpha\approx137$ is absorbed by
setting $e=1$ and the electron charge $-1$, giving
$\hat{\boldsymbol{\pi}} = \hat{\mathbf{p}} + \Avec$ as written.}
The canonical momentum is not gauge invariant and has no direct
physical significance;\cite{Sakurai-2nd_ed} the
\emph{mechanical} (kinetic) momentum
\begin{equation}
  \pihat \equiv \hat{\mathbf{p}} + \Avec(t) = -\imath\nabla + \Avec(t)
  \label{eq:pi-mech}
\end{equation}
is gauge invariant and represents the physical velocity of the
particle.\cite{Sakurai-2nd_ed,starace-1971}
In Gaussian basis set language, this is directly analogous to the
fact that only the kinetic energy integral $\langle\phi_\mu|T|
\phi_\nu\rangle$, not the kinetic canonical momentum itself, enters
the observable transition matrix elements.

Expanding the kinetic term in eq.~(\ref{eq:H-vg-2}):
\begin{align}
  \frac{1}{2}\pihat^2
  &= \frac{1}{2}
     \!\left(-\imath\nabla + \Avec\right)^{\!2}
  \nonumber \\
  &= \underbrace{-\tfrac{1}{2}\nabla^2}_{T_0}
   + \underbrace{\frac{\imath}{2}
     \left(\Avec\cdot\nabla + \nabla\cdot\Avec\right)
     }_{\text{paramagnetic coupling}}
   + \underbrace{\frac{1}{2}\Avec^2}_{\text{diamagnetic}},
  \label{eq:kinetic-expand}
\end{align}
where we use $\nabla\cdot\Avec = 0$ in the Coulomb gauge to
simplify $\Avec\cdot\nabla + \nabla\cdot\Avec = 2\Avec\cdot\nabla$.
In the linear response regime the $\Avec^2$ diamagnetic term is
dropped (it is second order in the applied field).

%--------------------------------------------------------------------
\subsubsection{Velocity Operator from the Heisenberg Equation of Motion}
\label{sec:velocity-eom}
%--------------------------------------------------------------------

The velocity operator is defined via the Heisenberg equation of
motion for the position operator:
\begin{equation}
  \hat{v}_\alpha
  \equiv \frac{1}{\imath}\left[\hat{r}_\alpha,\hat{H}\right].
  \label{eq:heisenberg-v}
\end{equation}
We evaluate eq.~(\ref{eq:heisenberg-v}) for each term in
$\hat{H} = \frac{1}{2}\pihat^2 + \Vloc + \Vnl$ separately.

\paragraph{Kinetic energy contribution.}
We compute $\frac{1}{\imath}[\hat{r}_\alpha, -\frac{1}{2}\nabla^2]$
by acting on an arbitrary test function $g(\mathbf{r})$:
\begin{align}
  \left[\hat{r}_\alpha,\nabla_\alpha^2\right]g
  &= r_\alpha\,\nabla_\alpha^2 g
   - \nabla_\alpha^2(r_\alpha g)
  \nonumber \\
  &= r_\alpha\,\nabla_\alpha^2 g
   - \nabla_\alpha\!\left(g + r_\alpha\nabla_\alpha g\right)
  \nonumber \\
  &= r_\alpha\,\nabla_\alpha^2 g
   - \nabla_\alpha g
   - \nabla_\alpha g
   - r_\alpha\nabla_\alpha^2 g
   = -2\nabla_\alpha g,
  \label{eq:comm-r-nabla2}
\end{align}
where we applied the product rule twice.
The same calculation holds for each Cartesian component, giving
$[\hat{r}_\alpha,\nabla^2] = -2\nabla_\alpha$ (where the right-hand
side involves only the $\alpha$-component of $\nabla$).
Therefore
\begin{equation}
  \frac{1}{\imath}
  \!\left[\hat{r}_\alpha,{-\tfrac{1}{2}\nabla^2}\right]
  = \frac{1}{\imath}\cdot(-\tfrac{1}{2})\cdot(-2\nabla_\alpha)
  = \frac{\nabla_\alpha}{\imath}
  = -\imath\nabla_\alpha = \hat{p}_\alpha.
  \label{eq:comm-r-T}
\end{equation}
Including the full mechanical momentum:
$\frac{1}{\imath}[\hat{r}_\alpha, \frac{1}{2}\pihat^2]
= \hat{p}_\alpha + A_\alpha(t) = \hat{\pi}_\alpha$.

\paragraph{Local potential contribution.}
Because $\Vloc(\mathbf{r})$ is multiplicative (diagonal in position
space),
\begin{equation}
  \left[\hat{r}_\alpha,\Vloc(\mathbf{r})\right]g
  = r_\alpha\,\Vloc g - \Vloc\,r_\alpha g = 0.
  \label{eq:comm-r-Vloc}
\end{equation}
This is the same reason that, in a Gaussian basis, the velocity
matrix elements $\langle\mu|-\imath\nabla|\nu\rangle$ and the
length matrix elements $(\varepsilon_\mu - \varepsilon_\nu)
\langle\mu|\hat{r}|\nu\rangle$ are equal when $V$ is local: both
gauges give the same result.

\paragraph{Nonlocal pseudopotential contribution.}
The nonlocal pseudopotential $\Vnl$ is not diagonal in position
space; its action involves an integral over all space,
\begin{equation}
  \left(\Vnl\psi\right)(\mathbf{r})
  = \int \langle\mathbf{r}|\Vnl|\mathbf{r}'\rangle\,
    \psi(\mathbf{r}')\,\mathrm{d}^3r',
  \label{eq:Vnl-nonlocal-action}
\end{equation}
so $\hat{r}_\alpha$ and $\Vnl$ do not share a common eigenbasis
and their commutator is nonzero.
This is directly analogous to the situation in Gaussian basis sets
with effective core potentials (ECPs), where the ECP commutator
$[\hat{r},V_\mathrm{ECP}]$ also gives a nonzero
contribution.\cite{starace-1971}
The explicit form of $\frac{1}{\imath}[\hat{r}_\alpha,\Vnl]$ for
Kleinman-Bylander separable pseudopotentials is derived in
Appendix~\ref{app:gauge_nlpp}.

\paragraph{Full result.}
Collecting all contributions, the velocity operator is
\begin{equation}
  \boxed{
  \hat{v}_\alpha
  = \frac{1}{\imath}\left[\hat{r}_\alpha,\hat{H}\right]
  = \underbrace{\hat{\pi}_\alpha}_{\displaystyle
    = -\imath\nabla_\alpha + A_\alpha}
  \;+\;
  \underbrace{\frac{1}{\imath}
    \left[\hat{r}_\alpha,\Vnl\right]}_{\text{nonlocal PP correction}}.
  }
  \label{eq:v-full}
\end{equation}
For a purely local Hamiltonian the second term vanishes and
$\hat{v}_\alpha = \hat{\pi}_\alpha$.
The nonlocal correction was first identified by
Starace\cite{starace-1971,starace-1973} in the context of
photoionization from atoms described by model potentials, and is
discussed in detail by Mattiat and Luber\cite{Mattiat-Luber-2022}
for RT-TDDFT with nonlocal pseudopotentials.
Its physical origin is that the electron ``feels'' the nonlocal
potential at points $\mathbf{r}'$ away from its instantaneous
position $\mathbf{r}$, so its velocity cannot be expressed through
the local gradient alone.

%--------------------------------------------------------------------
\subsubsection{Current Density from the Continuity Equation}
\label{sec:current-continuity}
%--------------------------------------------------------------------

We now derive the same current density by an independent route
through the continuity equation, which provides a consistency check
and clarifies the physical role of each contribution.

The probability density is $\rho(\mathbf{r},t) = \psi^*\psi$.
Differentiating with respect to time and substituting the TDKS
equation $\imath\partial_t\psi = \hat{H}\psi$ and its Hermitian
conjugate $-\imath\partial_t\psi^* = \hat{H}\psi^*$:
\begin{align}
  \frac{\partial\rho}{\partial t}
  &= \dot\psi^*\psi + \psi^*\dot\psi
   = \imath\!\left[(\hat{H}\psi^*)\psi
                  - \psi^*(\hat{H}\psi)\right].
  \label{eq:drho-H}
\end{align}
We evaluate the right-hand side term by term.

\paragraph{Local potential: vanishes.}
$\imath[(\Vloc\psi^*)\psi - \psi^*(\Vloc\psi)]
= \imath\Vloc(\psi^*\psi - \psi^*\psi) = 0$.

\paragraph{Kinetic energy: paramagnetic and diamagnetic currents.}
Using the Coulomb gauge $\nabla\cdot\Avec = 0$ and abbreviating
$\hat{D} \equiv -\imath\nabla + \Avec$:
\begin{align}
  \imath\!\left[(\hat{T}\psi^*)\psi - \psi^*(\hat{T}\psi)\right]
  &= \frac{\imath}{2}\!\left[
     \overline{(\hat{D}^2\psi)}\cdot\psi
     - \psi^*\cdot(\hat{D}^2\psi)
     \right]
  \nonumber \\
  &= -\nabla\cdot\!\underbrace{\left[
     \frac{1}{2}\operatorname{Re}\!\left(
     \psi^*\hat{D}\psi\right)\right]}_{\displaystyle
     = \Jpara + \Jdia},
  \label{eq:kin-continuity}
\end{align}
where the last equality follows from integration by parts applied
to $\nabla\cdot(\psi^*\hat{D}\psi)$.
Expanding $\hat{D} = -\imath\nabla + \Avec$ inside the
$\operatorname{Re}[\cdot]$:
\begin{align}
  \Jpara(\mathbf{r},t)
  &= \operatorname{Re}\!\left[\psi^*(-\imath\nabla)\psi\right],
  \label{eq:Jpara-def}
  \\
  \Jdia(\mathbf{r},t)
  &= \Avec(t)\,|\psi(\mathbf{r},t)|^2.
  \label{eq:Jdia-def}
\end{align}

\paragraph{Nonlocal pseudopotential: nonlocal current.}
The nonlocal term contributes
\begin{equation}
  \imath\!\left[(\Vnl\psi^*)\psi - \psi^*(\Vnl\psi)\right]
  = -\nabla\cdot\mathbf{J}_\mathrm{nl}(\mathbf{r},t),
  \label{eq:Jnl-continuity}
\end{equation}
where $\mathbf{J}_\mathrm{nl}$ is the nonlocal current density
whose explicit form in terms of the Kleinman-Bylander projectors
is derived in Appendix~\ref{app:Jnl-expression}.

\paragraph{Total current density.}
Comparing with the continuity equation
$\partial_t\rho + \nabla\cdot\mathbf{J} = 0$ and collecting
all contributions:
\begin{equation}
  \mathbf{J}(\mathbf{r},t)
  = \underbrace{
      \operatorname{Re}\!\left[\psi^*(-\imath\nabla)\psi\right]
    }_{\Jpara}
  + \underbrace{
      \Avec(t)|\psi|^2
    }_{\Jdia}
  + \underbrace{
      \operatorname{Re}\!\left[
        \psi^*\tfrac{1}{\imath}[\hat{\mathbf{r}},\Vnl]\psi
      \right]
    }_{\mathbf{J}_\mathrm{nl}}.
  \label{eq:J-total}
\end{equation}
This is fully consistent with eq.~(\ref{eq:v-full}): the three
contributions map one-to-one onto the mechanical momentum
$\hat{\pi}_\alpha$ (split into its $\hat{p}_\alpha$ and $A_\alpha$
parts) plus the nonlocal correction.

%--------------------------------------------------------------------
\subsubsection{Paramagnetic and Diamagnetic Contributions:
  Physical Roles}
\label{sec:para-dia}
%--------------------------------------------------------------------

The \emph{paramagnetic} current $\Jpara$ is independent of the
external field $\Avec$ and represents the intrinsic probability
current of the time-evolving quantum state.
It oscillates in response to the delta-kick perturbation and, after
Fourier transformation, yields the optical conductivity.

The \emph{diamagnetic} current $\Jdia = \Avec(t)|\psi|^2$ is
proportional to the applied vector potential.
In the perturbative kick scheme used in \texttt{RMG},%
\cite{yabana-bertsch-1996,jakowski-rmg-tddft-2025}
the vector potential is a delta-function impulse:
$\Avec(t) = \Avec_0\,\delta(t)$, with $\Avec(t) = 0$ for all
$t > 0$.
Consequently:
\begin{equation}
  \Jdia(\mathbf{r},t>0) = \Avec(t>0)\,|\psi|^2 = 0.
  \label{eq:Jdia-zero}
\end{equation}
The diamagnetic current vanishes identically during the simulation
after the kick.
This is not an approximation: within linear response theory the
complete first-order optical response is encoded in the oscillations
of $\Jpara(t>0)$ and $\mathbf{J}_\mathrm{nl}(t>0)$,%
\cite{yabana-bertsch-1996}
so the macroscopic current that is recorded and Fourier-transformed
is
\begin{equation}
  J_\alpha(t>0)
  = \int_\Omega\!
    \left[\Jpara(\mathbf{r},t)
    + \mathbf{J}_{\mathrm{nl}}(\mathbf{r},t)\right]_\alpha
    \mathrm{d}^3r,
  \label{eq:J-sim}
\end{equation}
where the integral is over the unit cell $\Omega$.

%--------------------------------------------------------------------
\subsubsection{Current Operator for Bloch States:
  the $\mathbf{k}$-Vector Correction}
\label{sec:current-bloch}
%--------------------------------------------------------------------

For a periodic crystal, the Kohn-Sham orbitals satisfy Bloch's
theorem.
We derive how the current operator is modified when expressed in
terms of the cell-periodic part $u_{n,\mathbf{k}}$.

A Bloch state at band index $n$ and crystal momentum $\mathbf{k}$
is written as
\begin{equation}
  \psi_{n,\mathbf{k}}(\mathbf{r},t)
  = e^{\imath\mathbf{k}\cdot\mathbf{r}}\,
    u_{n,\mathbf{k}}(\mathbf{r},t),
  \label{eq:bloch-def}
\end{equation}
where $u_{n,\mathbf{k}}(\mathbf{r}+\mathbf{R},t)
= u_{n,\mathbf{k}}(\mathbf{r},t)$ for all lattice vectors
$\mathbf{R}$.
Applying the canonical momentum operator $-\imath\nabla_\alpha$ to
the full Bloch state and using the product rule:
\begin{align}
  (-\imath\nabla_\alpha)\,\psi_{n,\mathbf{k}}
  &= (-\imath\nabla_\alpha)\!\left[
     e^{\imath\mathbf{k}\cdot\mathbf{r}}\,u_{n,\mathbf{k}}
     \right]
  \nonumber \\
  &= (-\imath)\!\left[
     \imath k_\alpha\,
     e^{\imath\mathbf{k}\cdot\mathbf{r}}\,u_{n,\mathbf{k}}
     + e^{\imath\mathbf{k}\cdot\mathbf{r}}\,
       \nabla_\alpha u_{n,\mathbf{k}}
     \right]
  \nonumber \\
  &= e^{\imath\mathbf{k}\cdot\mathbf{r}}
     \!\left[k_\alpha + (-\imath\nabla_\alpha)\right]
     u_{n,\mathbf{k}}.
  \label{eq:nabla-bloch}
\end{align}
The Bloch phase factor $e^{\imath\mathbf{k}\cdot\mathbf{r}}$
cancels between bra and ket in a matrix element over the unit cell,
so the matrix element of the canonical momentum between two Bloch
states at the same $\mathbf{k}$-point is
\begin{align}
  p^\alpha_{ab}(\mathbf{k})
  &\equiv \langle\psi_{a,\mathbf{k}}|(-\imath\nabla_\alpha)|
    \psi_{b,\mathbf{k}}\rangle
  \nonumber \\
  &= \int_\Omega\!
     e^{-\imath\mathbf{k}\cdot\mathbf{r}}\,u_{a,\mathbf{k}}^*\cdot
     e^{\imath\mathbf{k}\cdot\mathbf{r}}
     \!\left(-\imath\nabla_\alpha + k_\alpha\right)
     u_{b,\mathbf{k}}\,\mathrm{d}^3r
  \nonumber \\
  &= \langle u_{a,\mathbf{k}}|
     \!\left(-\imath\nabla_\alpha + k_\alpha\right)
     |u_{b,\mathbf{k}}\rangle_\Omega.
  \label{eq:para-bloch-me}
\end{align}
The $k_\alpha$ shift arises purely from the Bloch factor and is
absent for $\Gamma$-point ($\mathbf{k} = \mathbf{0}$) calculations.
For finite $\mathbf{k}$-point sampling it must be included to obtain
the correct optical response; its omission would give wrong
intensities proportional to $|\mathbf{k}|^2$ at each $\mathbf{k}$-point.

The full current density summed over all occupied Bloch states is
\begin{equation}
  \mathbf{J}(\mathbf{r},t)
  = \frac{1}{N_k}\!\sum_{n,\mathbf{k}}^{\mathrm{occ}}
    f_{n,\mathbf{k}}\,
    \operatorname{Re}\!\left[
      u_{n,\mathbf{k}}^*\!\left(
      -\imath\nabla + \mathbf{k}
      + \frac{1}{\imath}[\hat{\mathbf{r}},\Vnl]
      \right)u_{n,\mathbf{k}}
    \right],
  \label{eq:J-bloch-total}
\end{equation}
where $N_k$ is the number of $\mathbf{k}$-points and $f_{n,\mathbf{k}}$
are the occupation numbers.
Equation~(\ref{eq:J-bloch-total}) reproduces eq.~(8) of Yabana and
Bertsch\cite{yabana-bertsch-1996} and eq.~(26) of Mattiat and
Luber,\cite{Mattiat-Luber-2022} specialized to the Coulomb gauge
and the velocity gauge with a spatially uniform $\Avec(t)$.

%--------------------------------------------------------------------
\subsubsection{Current Operator Matrix and Density Matrix Formulation}
\label{sec:current-dm-sec}
%--------------------------------------------------------------------

In the density matrix formulation the macroscopic current is
evaluated by contracting the propagated density matrix
$\mathbf{P}_\mathbf{k}(t)$ (Section~\ref{sec:propagation}) with
a fixed current operator matrix that is assembled once before the
time loop.
The current operator matrix element in the Kohn-Sham orbital basis
at $\mathbf{k}$-point $\mathbf{k}$ is
\begin{equation}
  \boxed{
  \Jmat{\alpha}_{ab}(\mathbf{k})
  = \underbrace{
      \langle\psi_{a,\mathbf{k}}|
      (-\imath\nabla_\alpha + k_\alpha)
      |\psi_{b,\mathbf{k}}\rangle
    }_{\displaystyle\texttt{VecPHmatrix()}}
  +
  \underbrace{
      \langle\psi_{a,\mathbf{k}}|
      \tfrac{1}{\imath}[\hat{r}_\alpha,\Vnl]
      |\psi_{b,\mathbf{k}}\rangle
    }_{\displaystyle\texttt{CurrentNlpp()}},
  }
  \label{eq:Jmat-full}
\end{equation}
where the code labels refer to \texttt{RMG} routines described in
Section~\ref{sec:implementation}.
The macroscopic $\alpha$-component of the current is then
\begin{equation}
  J_\alpha(t)
  = \sum_{\mathbf{k}} w_\mathbf{k}\,
    \operatorname{Re}\operatorname{Tr}\!\left[
      \mathbf{P}_\mathbf{k}(t)\cdot
      \boldsymbol{\mathcal{J}}^\alpha_\mathbf{k}
    \right],
  \label{eq:J-dm}
\end{equation}
where $w_\mathbf{k}$ are the $\mathbf{k}$-point weights
(with $\sum_\mathbf{k} w_\mathbf{k} = 1$) and the trace runs over
the active orbital subspace.
In the code, eq.~(\ref{eq:J-dm}) is evaluated at each time step by
a single BLAS complex dot product:
\begin{verbatim}
J[alpha] += Re(zdotc(Norb^2, Pn0, 1, Palpha_matrix, 1)) * kweight;
\end{verbatim}
where \texttt{Palpha\_matrix} contains the stored matrix
$\Jmat{\alpha}_{ab}$ (built by \texttt{VecPHmatrix} +
\texttt{CurrentNlpp} before the time loop) and \texttt{Pn0}
contains the current density matrix elements $P_{ab}(t)$.

%--------------------------------------------------------------------
\subsubsection{Hypervirial Relation and Gauge Invariance}
\label{sec:hypervirial}
%--------------------------------------------------------------------

The hypervirial relation is a consistency condition that connects the
velocity (current) operator in the velocity gauge to the dipole
matrix elements in the length gauge.
It is the periodic-system analog of the relation between oscillator
strengths in the length and velocity forms.

In the Heisenberg picture, the Ehrenfest theorem applied to
$\hat{r}_\alpha$ gives:
\begin{equation}
  \langle\psi_a|\hat{v}_\alpha|\psi_b\rangle
  = \langle\psi_a|\tfrac{1}{\imath}[\hat{r}_\alpha,\hat{H}]
    |\psi_b\rangle
  = (\varepsilon_b - \varepsilon_a)
    \langle\psi_a|\hat{r}_\alpha|\psi_b\rangle,
  \label{eq:ehrenfest}
\end{equation}
where $\varepsilon_a$, $\varepsilon_b$ are the ground-state
Kohn-Sham eigenvalues.
For a Hamiltonian with only local potentials, where $\hat{v}_\alpha
= -\imath\nabla_\alpha$, this gives the standard hypervirial relation
\begin{equation}
  \langle\psi_a|(-\imath\nabla_\alpha)|\psi_b\rangle
  = (\varepsilon_b - \varepsilon_a)
    \langle\psi_a|\hat{r}_\alpha|\psi_b\rangle,
  \label{eq:hypervirial-local}
\end{equation}
which states that the velocity gauge matrix element (left side)
equals the length gauge matrix element (right side): both gauges
are equivalent.

When the Hamiltonian contains a nonlocal pseudopotential, the
velocity operator gains the correction
$\frac{1}{\imath}[\hat{r}_\alpha,\Vnl]$
(eq.~\ref{eq:v-full}).
Inserting into eq.~(\ref{eq:ehrenfest}), and rearranging:
\begin{equation}
  \boxed{
  \langle\psi_a|(-\imath\nabla_\alpha)|\psi_b\rangle
  = (\varepsilon_b - \varepsilon_a)
    \langle\psi_a|\hat{r}_\alpha|\psi_b\rangle
  - \langle\psi_a|\tfrac{1}{\imath}
    [\hat{r}_\alpha,\Vnl]|\psi_b\rangle.
  }
  \label{eq:hypervirial-nlpp}
\end{equation}
Equation~(\ref{eq:hypervirial-nlpp}) is the nonlocal-PP-corrected
hypervirial relation, first given by Starace\cite{starace-1971} and
discussed for RT-TDDFT by Mattiat and
Luber.\cite{Mattiat-Luber-2022}
Its physical meaning is that omitting the commutator term from the
velocity operator would give different results in the velocity and
length gauges — that is, wrong spectra.
Including the correction, and only then, restores the gauge
equivalence and ensures that the computed optical conductivity is
independent of the chosen gauge.
Equation~(\ref{eq:hypervirial-nlpp}) also provides a practical
numerical diagnostic: the left-hand side is computed by
\texttt{VecPHmatrix()}, while the right-hand side can be computed
independently from the Kohn-Sham eigenvalues and the dipole matrix;
agreement between the two, up to the grid discretization error,
validates the implementation.

%--------------------------------------------------------------------
\subsubsection{Sign Convention and Factor of $\imath$ in the Code}
\label{sec:imath-convention}
%--------------------------------------------------------------------

We track the factors of $\imath$ explicitly to verify that the code
implements eq.~(\ref{eq:J-dm}) with the correct sign.

The physical current matrix element
$\Jmat{\alpha}_{ab} = \langle\psi_a|(-\imath\nabla_\alpha + k_\alpha
+ \frac{1}{\imath}[\hat{r}_\alpha,\Vnl])|\psi_b\rangle$
is complex-valued in general.
In \texttt{VecPHmatrix()}, the gradient matrix is computed and stored
multiplied by the prefactor \texttt{I\_t} $= \imath$:
\begin{equation}
  \texttt{Pxmatrix}[ab]
  \leftarrow
  \imath\cdot\langle\psi_a|(-\imath\nabla_\alpha)|\psi_b\rangle
  + \imath\cdot k_\alpha\langle\psi_a|\psi_b\rangle
  = \langle\psi_a|(\nabla_\alpha - \imath k_\alpha)|\psi_b\rangle.
  \label{eq:stored-para}
\end{equation}
That is, \texttt{Pxmatrix} stores $\imath\Jmat{\alpha}_{ab}$ from
the paramagnetic part.
\texttt{CurrentNlpp()} adds the nonlocal contribution with the same
prefactor:
\begin{equation}
  \texttt{Pxmatrix}[ab]
  \mathrel{+}=
  \imath\cdot\langle\psi_a|
  \tfrac{1}{\imath}[\hat{r}_\alpha,\Vnl]|\psi_b\rangle
  = \langle\psi_a|[\hat{r}_\alpha,\Vnl]|\psi_b\rangle.
  \label{eq:stored-nlpp}
\end{equation}
Thus \texttt{Pxmatrix} $= \imath\cdot\Jmat{\alpha}_{ab}$ in total.
The macroscopic current is extracted as
\begin{align}
  J_\alpha
  &= \operatorname{Re}\operatorname{Tr}
     \!\left[\mathbf{P}\cdot\texttt{Pxmatrix}\right]
   = \operatorname{Re}\operatorname{Tr}
     \!\left[\mathbf{P}\cdot(\imath\boldsymbol{\mathcal{J}}^\alpha)\right].
  \label{eq:J-from-code}
\end{align}
For a Hermitian density matrix $\mathbf{P} = \mathbf{P}^\dag$
and an anti-Hermitian current operator matrix
$(\boldsymbol{\mathcal{J}}^\alpha)^\dag = -\boldsymbol{\mathcal{J}}^\alpha$
(which follows from the anti-Hermitian nature of $-\imath\nabla$
and of $\frac{1}{\imath}[\hat{r},\Vnl]$%
\footnote{Both $-\imath\nabla$ and $\frac{1}{\imath}[\hat{r},\Vnl]$
are anti-Hermitian operators when the projectors $\beta^{lm}_I$
are real-valued, which is the case for norm-conserving
pseudopotentials in \texttt{RMG}.}),
the trace satisfies
$\operatorname{Tr}[\mathbf{P}\boldsymbol{\mathcal{J}}^\alpha]
= -\overline{\operatorname{Tr}[\mathbf{P}\boldsymbol{\mathcal{J}}^\alpha]}$,
i.e., it is purely imaginary.
Therefore
\begin{equation}
  J_\alpha
  = \operatorname{Re}\!\left[\imath\operatorname{Tr}
    (\mathbf{P}\boldsymbol{\mathcal{J}}^\alpha)\right]
  = -\operatorname{Im}\operatorname{Tr}
    \!\left[\mathbf{P}\boldsymbol{\mathcal{J}}^\alpha\right]
  = \operatorname{Re}\operatorname{Tr}
    \!\left[\mathbf{P}\cdot\imath\boldsymbol{\mathcal{J}}^\alpha\right],
  \label{eq:J-sign-check}
\end{equation}
which is consistent with eq.~(\ref{eq:J-dm}) and confirms that the
sign convention is correct.
The uniform prefactor \texttt{I\_t} applied in both
\texttt{VecPHmatrix()} and \texttt{CurrentNlpp()} ensures that the
paramagnetic and nonlocal contributions enter the current with
the same relative sign, as required by eq.~(\ref{eq:Jmat-full}).



 replaces the placeholder in sec_theory.tex.
%%
%% CONVENTIONS:
%%   Atomic units: hbar = m_e = e = 4*pi*eps_0 = 1
%%   Electron charge = -1  (e > 0 is the elementary charge magnitude)
%%   In atomic units the Hamiltonian is
%%       H = (1/2)(-i*nabla + A)^2 + V_loc + V_nl
%%   Imaginary unit: \imath  (dotless i, consistent with paper 1)
%%   Sign check: NEAT-notes-RMG has a typo (-1/2m kinetic energy);
%%               NEAT-project-notes has the correct +1/2m sign — used here.
%%
%% NOTATION MACROS (all in notation.tex):
%%   \phat   = \hat{\mathbf{p}}  canonical momentum
%%   \pihat  = \hat{\boldsymbol{\pi}}  mechanical momentum
%%   \vhat   = \hat{\mathbf{v}}  velocity operator
%%   \rhat   = \hat{\mathbf{r}}  position
%%   \Vnl    = V^{\mathrm{nl}}
%%   \Vloc   = V^{\mathrm{loc}}
%%   \Avec   = \mathbf{A}
%%   \Jpara  = \mathbf{J}_{\mathrm{para}}
%%   \Jdia   = \mathbf{J}_{\mathrm{dia}}
%%   \Jmat{x}= \mathcal{J}^{x}
%%%%%%%%%%%%%%%%%%%%%%%%%%%%%%%%%%%%%%%%%%%%%%%%%%%%%%%%%%%%%%%%%%%%%%%%%%%%%%%

\subsection{Current Operator and Nonlocal Pseudopotential Contributions}
\label{sec:current}

The central observable in periodic RT-TDDFT is the macroscopic
electronic current density $J_\alpha(t)$, from which the optical
conductivity $\sigma_{\alpha\beta}(\omega)$ and dielectric function
$\varepsilon_{\alpha\beta}(\omega)$ are obtained by Fourier analysis
(Section~\ref{sec:postprocessing}).
For periodic systems the electric dipole moment
$\langle\hat{\mathbf{r}}\rangle$ is ill-defined because the position
operator is incompatible with periodic boundary
conditions,\cite{resta-1998,king-smith-vanderbilt-1993} and the
current is the natural replacement.

We derive the current density operator in four steps: (i) we
distinguish canonical from mechanical momentum and define the velocity
operator via the Heisenberg equation of motion, identifying the role
of nonlocal pseudopotentials; (ii) we rederive the same result from
the continuity equation, decomposing the current into paramagnetic,
diamagnetic, and nonlocal pseudopotential contributions; (iii) we
specialize to Bloch states and derive the $\mathbf{k}$-vector
correction explicitly from the Bloch ansatz; (iv) we express the
macroscopic current in density matrix form and verify the sign
convention in the code.
All equations are in atomic units ($\hbar = m_e = e = 4\pi\varepsilon_0
= 1$) with the electron charge equal to $-1$.

%--------------------------------------------------------------------
\subsubsection{Canonical and Mechanical Momentum}
\label{sec:canon-mech}
%--------------------------------------------------------------------

The Hamiltonian of the Kohn-Sham system in the presence of a
spatially uniform, time-dependent vector potential $\Avec(t)$
(long-wavelength approximation, see Section~\ref{sec:EM-coupling})
is
\begin{equation}
  \hat{H}(t) = \frac{1}{2}
    \!\left(\hat{\mathbf{p}} + \Avec(t)\right)^{\!2}
    + \Vloc(\mathbf{r}) + \Vnl,
  \label{eq:H-vg-2}
\end{equation}
where $\hat{\mathbf{p}} = -\imath\nabla$ is the \emph{canonical}
momentum operator and the sign on $\Avec$ reflects that in atomic
units with electron charge $-1$ the coupling is
$q\Avec = (+1)\cdot(-1)\cdot\Avec = -\Avec$, giving mechanical
momentum $\pihat = \hat{\mathbf{p}} + \Avec$.\footnote{In
Gaussian units the coupling is $q\Avec/c = -e\Avec/c$ for an
electron; in atomic units $c=1/\alpha\approx137$ is absorbed by
setting $e=1$ and the electron charge $-1$, giving
$\hat{\boldsymbol{\pi}} = \hat{\mathbf{p}} + \Avec$ as written.}
The canonical momentum is not gauge invariant and has no direct
physical significance;\cite{Sakurai-2nd_ed} the
\emph{mechanical} (kinetic) momentum
\begin{equation}
  \pihat \equiv \hat{\mathbf{p}} + \Avec(t) = -\imath\nabla + \Avec(t)
  \label{eq:pi-mech}
\end{equation}
is gauge invariant and represents the physical velocity of the
particle.\cite{Sakurai-2nd_ed,starace-1971}
In Gaussian basis set language, this is directly analogous to the
fact that only the kinetic energy integral $\langle\phi_\mu|T|
\phi_\nu\rangle$, not the kinetic canonical momentum itself, enters
the observable transition matrix elements.

Expanding the kinetic term in eq.~(\ref{eq:H-vg-2}):
\begin{align}
  \frac{1}{2}\pihat^2
  &= \frac{1}{2}
     \!\left(-\imath\nabla + \Avec\right)^{\!2}
  \nonumber \\
  &= \underbrace{-\tfrac{1}{2}\nabla^2}_{T_0}
   + \underbrace{\frac{\imath}{2}
     \left(\Avec\cdot\nabla + \nabla\cdot\Avec\right)
     }_{\text{paramagnetic coupling}}
   + \underbrace{\frac{1}{2}\Avec^2}_{\text{diamagnetic}},
  \label{eq:kinetic-expand}
\end{align}
where we use $\nabla\cdot\Avec = 0$ in the Coulomb gauge to
simplify $\Avec\cdot\nabla + \nabla\cdot\Avec = 2\Avec\cdot\nabla$.
In the linear response regime the $\Avec^2$ diamagnetic term is
dropped (it is second order in the applied field).

%--------------------------------------------------------------------
\subsubsection{Velocity Operator from the Heisenberg Equation of Motion}
\label{sec:velocity-eom}
%--------------------------------------------------------------------

The velocity operator is defined via the Heisenberg equation of
motion for the position operator:
\begin{equation}
  \hat{v}_\alpha
  \equiv \frac{1}{\imath}\left[\hat{r}_\alpha,\hat{H}\right].
  \label{eq:heisenberg-v}
\end{equation}
We evaluate eq.~(\ref{eq:heisenberg-v}) for each term in
$\hat{H} = \frac{1}{2}\pihat^2 + \Vloc + \Vnl$ separately.

\paragraph{Kinetic energy contribution.}
We compute $\frac{1}{\imath}[\hat{r}_\alpha, -\frac{1}{2}\nabla^2]$
by acting on an arbitrary test function $g(\mathbf{r})$:
\begin{align}
  \left[\hat{r}_\alpha,\nabla_\alpha^2\right]g
  &= r_\alpha\,\nabla_\alpha^2 g
   - \nabla_\alpha^2(r_\alpha g)
  \nonumber \\
  &= r_\alpha\,\nabla_\alpha^2 g
   - \nabla_\alpha\!\left(g + r_\alpha\nabla_\alpha g\right)
  \nonumber \\
  &= r_\alpha\,\nabla_\alpha^2 g
   - \nabla_\alpha g
   - \nabla_\alpha g
   - r_\alpha\nabla_\alpha^2 g
   = -2\nabla_\alpha g,
  \label{eq:comm-r-nabla2}
\end{align}
where we applied the product rule twice.
The same calculation holds for each Cartesian component, giving
$[\hat{r}_\alpha,\nabla^2] = -2\nabla_\alpha$ (where the right-hand
side involves only the $\alpha$-component of $\nabla$).
Therefore
\begin{equation}
  \frac{1}{\imath}
  \!\left[\hat{r}_\alpha,{-\tfrac{1}{2}\nabla^2}\right]
  = \frac{1}{\imath}\cdot(-\tfrac{1}{2})\cdot(-2\nabla_\alpha)
  = \frac{\nabla_\alpha}{\imath}
  = -\imath\nabla_\alpha = \hat{p}_\alpha.
  \label{eq:comm-r-T}
\end{equation}
Including the full mechanical momentum:
$\frac{1}{\imath}[\hat{r}_\alpha, \frac{1}{2}\pihat^2]
= \hat{p}_\alpha + A_\alpha(t) = \hat{\pi}_\alpha$.

\paragraph{Local potential contribution.}
Because $\Vloc(\mathbf{r})$ is multiplicative (diagonal in position
space),
\begin{equation}
  \left[\hat{r}_\alpha,\Vloc(\mathbf{r})\right]g
  = r_\alpha\,\Vloc g - \Vloc\,r_\alpha g = 0.
  \label{eq:comm-r-Vloc}
\end{equation}
This is the same reason that, in a Gaussian basis, the velocity
matrix elements $\langle\mu|-\imath\nabla|\nu\rangle$ and the
length matrix elements $(\varepsilon_\mu - \varepsilon_\nu)
\langle\mu|\hat{r}|\nu\rangle$ are equal when $V$ is local: both
gauges give the same result.

\paragraph{Nonlocal pseudopotential contribution.}
The nonlocal pseudopotential $\Vnl$ is not diagonal in position
space; its action involves an integral over all space,
\begin{equation}
  \left(\Vnl\psi\right)(\mathbf{r})
  = \int \langle\mathbf{r}|\Vnl|\mathbf{r}'\rangle\,
    \psi(\mathbf{r}')\,\mathrm{d}^3r',
  \label{eq:Vnl-nonlocal-action}
\end{equation}
so $\hat{r}_\alpha$ and $\Vnl$ do not share a common eigenbasis
and their commutator is nonzero.
This is directly analogous to the situation in Gaussian basis sets
with effective core potentials (ECPs), where the ECP commutator
$[\hat{r},V_\mathrm{ECP}]$ also gives a nonzero
contribution.\cite{starace-1971}
The explicit form of $\frac{1}{\imath}[\hat{r}_\alpha,\Vnl]$ for
Kleinman-Bylander separable pseudopotentials is derived in
Appendix~\ref{app:gauge_nlpp}.

\paragraph{Full result.}
Collecting all contributions, the velocity operator is
\begin{equation}
  \boxed{
  \hat{v}_\alpha
  = \frac{1}{\imath}\left[\hat{r}_\alpha,\hat{H}\right]
  = \underbrace{\hat{\pi}_\alpha}_{\displaystyle
    = -\imath\nabla_\alpha + A_\alpha}
  \;+\;
  \underbrace{\frac{1}{\imath}
    \left[\hat{r}_\alpha,\Vnl\right]}_{\text{nonlocal PP correction}}.
  }
  \label{eq:v-full}
\end{equation}
For a purely local Hamiltonian the second term vanishes and
$\hat{v}_\alpha = \hat{\pi}_\alpha$.
The nonlocal correction was first identified by
Starace\cite{starace-1971,starace-1973} in the context of
photoionization from atoms described by model potentials, and is
discussed in detail by Mattiat and Luber\cite{Mattiat-Luber-2022}
for RT-TDDFT with nonlocal pseudopotentials.
Its physical origin is that the electron ``feels'' the nonlocal
potential at points $\mathbf{r}'$ away from its instantaneous
position $\mathbf{r}$, so its velocity cannot be expressed through
the local gradient alone.

%--------------------------------------------------------------------
\subsubsection{Current Density from the Continuity Equation}
\label{sec:current-continuity}
%--------------------------------------------------------------------

We now derive the same current density by an independent route
through the continuity equation, which provides a consistency check
and clarifies the physical role of each contribution.

The probability density is $\rho(\mathbf{r},t) = \psi^*\psi$.
Differentiating with respect to time and substituting the TDKS
equation $\imath\partial_t\psi = \hat{H}\psi$ and its Hermitian
conjugate $-\imath\partial_t\psi^* = \hat{H}\psi^*$:
\begin{align}
  \frac{\partial\rho}{\partial t}
  &= \dot\psi^*\psi + \psi^*\dot\psi
   = \imath\!\left[(\hat{H}\psi^*)\psi
                  - \psi^*(\hat{H}\psi)\right].
  \label{eq:drho-H}
\end{align}
We evaluate the right-hand side term by term.

\paragraph{Local potential: vanishes.}
$\imath[(\Vloc\psi^*)\psi - \psi^*(\Vloc\psi)]
= \imath\Vloc(\psi^*\psi - \psi^*\psi) = 0$.

\paragraph{Kinetic energy: paramagnetic and diamagnetic currents.}
Using the Coulomb gauge $\nabla\cdot\Avec = 0$ and abbreviating
$\hat{D} \equiv -\imath\nabla + \Avec$:
\begin{align}
  \imath\!\left[(\hat{T}\psi^*)\psi - \psi^*(\hat{T}\psi)\right]
  &= \frac{\imath}{2}\!\left[
     \overline{(\hat{D}^2\psi)}\cdot\psi
     - \psi^*\cdot(\hat{D}^2\psi)
     \right]
  \nonumber \\
  &= -\nabla\cdot\!\underbrace{\left[
     \frac{1}{2}\operatorname{Re}\!\left(
     \psi^*\hat{D}\psi\right)\right]}_{\displaystyle
     = \Jpara + \Jdia},
  \label{eq:kin-continuity}
\end{align}
where the last equality follows from integration by parts applied
to $\nabla\cdot(\psi^*\hat{D}\psi)$.
Expanding $\hat{D} = -\imath\nabla + \Avec$ inside the
$\operatorname{Re}[\cdot]$:
\begin{align}
  \Jpara(\mathbf{r},t)
  &= \operatorname{Re}\!\left[\psi^*(-\imath\nabla)\psi\right],
  \label{eq:Jpara-def}
  \\
  \Jdia(\mathbf{r},t)
  &= \Avec(t)\,|\psi(\mathbf{r},t)|^2.
  \label{eq:Jdia-def}
\end{align}

\paragraph{Nonlocal pseudopotential: nonlocal current.}
The nonlocal term contributes
\begin{equation}
  \imath\!\left[(\Vnl\psi^*)\psi - \psi^*(\Vnl\psi)\right]
  = -\nabla\cdot\mathbf{J}_\mathrm{nl}(\mathbf{r},t),
  \label{eq:Jnl-continuity}
\end{equation}
where $\mathbf{J}_\mathrm{nl}$ is the nonlocal current density
whose explicit form in terms of the Kleinman-Bylander projectors
is derived in Appendix~\ref{app:Jnl-expression}.

\paragraph{Total current density.}
Comparing with the continuity equation
$\partial_t\rho + \nabla\cdot\mathbf{J} = 0$ and collecting
all contributions:
\begin{equation}
  \mathbf{J}(\mathbf{r},t)
  = \underbrace{
      \operatorname{Re}\!\left[\psi^*(-\imath\nabla)\psi\right]
    }_{\Jpara}
  + \underbrace{
      \Avec(t)|\psi|^2
    }_{\Jdia}
  + \underbrace{
      \operatorname{Re}\!\left[
        \psi^*\tfrac{1}{\imath}[\hat{\mathbf{r}},\Vnl]\psi
      \right]
    }_{\mathbf{J}_\mathrm{nl}}.
  \label{eq:J-total}
\end{equation}
This is fully consistent with eq.~(\ref{eq:v-full}): the three
contributions map one-to-one onto the mechanical momentum
$\hat{\pi}_\alpha$ (split into its $\hat{p}_\alpha$ and $A_\alpha$
parts) plus the nonlocal correction.

%--------------------------------------------------------------------
\subsubsection{Paramagnetic and Diamagnetic Contributions:
  Physical Roles}
\label{sec:para-dia}
%--------------------------------------------------------------------

The \emph{paramagnetic} current $\Jpara$ is independent of the
external field $\Avec$ and represents the intrinsic probability
current of the time-evolving quantum state.
It oscillates in response to the delta-kick perturbation and, after
Fourier transformation, yields the optical conductivity.

The \emph{diamagnetic} current $\Jdia = \Avec(t)|\psi|^2$ is
proportional to the applied vector potential.
In the perturbative kick scheme used in \texttt{RMG},%
\cite{yabana-bertsch-1996,jakowski-rmg-tddft-2025}
the vector potential is a delta-function impulse:
$\Avec(t) = \Avec_0\,\delta(t)$, with $\Avec(t) = 0$ for all
$t > 0$.
Consequently:
\begin{equation}
  \Jdia(\mathbf{r},t>0) = \Avec(t>0)\,|\psi|^2 = 0.
  \label{eq:Jdia-zero}
\end{equation}
The diamagnetic current vanishes identically during the simulation
after the kick.
This is not an approximation: within linear response theory the
complete first-order optical response is encoded in the oscillations
of $\Jpara(t>0)$ and $\mathbf{J}_\mathrm{nl}(t>0)$,%
\cite{yabana-bertsch-1996}
so the macroscopic current that is recorded and Fourier-transformed
is
\begin{equation}
  J_\alpha(t>0)
  = \int_\Omega\!
    \left[\Jpara(\mathbf{r},t)
    + \mathbf{J}_{\mathrm{nl}}(\mathbf{r},t)\right]_\alpha
    \mathrm{d}^3r,
  \label{eq:J-sim}
\end{equation}
where the integral is over the unit cell $\Omega$.

%--------------------------------------------------------------------
\subsubsection{Current Operator for Bloch States:
  the $\mathbf{k}$-Vector Correction}
\label{sec:current-bloch}
%--------------------------------------------------------------------

For a periodic crystal, the Kohn-Sham orbitals satisfy Bloch's
theorem.
We derive how the current operator is modified when expressed in
terms of the cell-periodic part $u_{n,\mathbf{k}}$.

A Bloch state at band index $n$ and crystal momentum $\mathbf{k}$
is written as
\begin{equation}
  \psi_{n,\mathbf{k}}(\mathbf{r},t)
  = e^{\imath\mathbf{k}\cdot\mathbf{r}}\,
    u_{n,\mathbf{k}}(\mathbf{r},t),
  \label{eq:bloch-def}
\end{equation}
where $u_{n,\mathbf{k}}(\mathbf{r}+\mathbf{R},t)
= u_{n,\mathbf{k}}(\mathbf{r},t)$ for all lattice vectors
$\mathbf{R}$.
Applying the canonical momentum operator $-\imath\nabla_\alpha$ to
the full Bloch state and using the product rule:
\begin{align}
  (-\imath\nabla_\alpha)\,\psi_{n,\mathbf{k}}
  &= (-\imath\nabla_\alpha)\!\left[
     e^{\imath\mathbf{k}\cdot\mathbf{r}}\,u_{n,\mathbf{k}}
     \right]
  \nonumber \\
  &= (-\imath)\!\left[
     \imath k_\alpha\,
     e^{\imath\mathbf{k}\cdot\mathbf{r}}\,u_{n,\mathbf{k}}
     + e^{\imath\mathbf{k}\cdot\mathbf{r}}\,
       \nabla_\alpha u_{n,\mathbf{k}}
     \right]
  \nonumber \\
  &= e^{\imath\mathbf{k}\cdot\mathbf{r}}
     \!\left[k_\alpha + (-\imath\nabla_\alpha)\right]
     u_{n,\mathbf{k}}.
  \label{eq:nabla-bloch}
\end{align}
The Bloch phase factor $e^{\imath\mathbf{k}\cdot\mathbf{r}}$
cancels between bra and ket in a matrix element over the unit cell,
so the matrix element of the canonical momentum between two Bloch
states at the same $\mathbf{k}$-point is
\begin{align}
  p^\alpha_{ab}(\mathbf{k})
  &\equiv \langle\psi_{a,\mathbf{k}}|(-\imath\nabla_\alpha)|
    \psi_{b,\mathbf{k}}\rangle
  \nonumber \\
  &= \int_\Omega\!
     e^{-\imath\mathbf{k}\cdot\mathbf{r}}\,u_{a,\mathbf{k}}^*\cdot
     e^{\imath\mathbf{k}\cdot\mathbf{r}}
     \!\left(-\imath\nabla_\alpha + k_\alpha\right)
     u_{b,\mathbf{k}}\,\mathrm{d}^3r
  \nonumber \\
  &= \langle u_{a,\mathbf{k}}|
     \!\left(-\imath\nabla_\alpha + k_\alpha\right)
     |u_{b,\mathbf{k}}\rangle_\Omega.
  \label{eq:para-bloch-me}
\end{align}
The $k_\alpha$ shift arises purely from the Bloch factor and is
absent for $\Gamma$-point ($\mathbf{k} = \mathbf{0}$) calculations.
For finite $\mathbf{k}$-point sampling it must be included to obtain
the correct optical response; its omission would give wrong
intensities proportional to $|\mathbf{k}|^2$ at each $\mathbf{k}$-point.

The full current density summed over all occupied Bloch states is
\begin{equation}
  \mathbf{J}(\mathbf{r},t)
  = \frac{1}{N_k}\!\sum_{n,\mathbf{k}}^{\mathrm{occ}}
    f_{n,\mathbf{k}}\,
    \operatorname{Re}\!\left[
      u_{n,\mathbf{k}}^*\!\left(
      -\imath\nabla + \mathbf{k}
      + \frac{1}{\imath}[\hat{\mathbf{r}},\Vnl]
      \right)u_{n,\mathbf{k}}
    \right],
  \label{eq:J-bloch-total}
\end{equation}
where $N_k$ is the number of $\mathbf{k}$-points and $f_{n,\mathbf{k}}$
are the occupation numbers.
Equation~(\ref{eq:J-bloch-total}) reproduces eq.~(8) of Yabana and
Bertsch\cite{yabana-bertsch-1996} and eq.~(26) of Mattiat and
Luber,\cite{Mattiat-Luber-2022} specialized to the Coulomb gauge
and the velocity gauge with a spatially uniform $\Avec(t)$.

%--------------------------------------------------------------------
\subsubsection{Current Operator Matrix and Density Matrix Formulation}
\label{sec:current-dm-sec}
%--------------------------------------------------------------------

In the density matrix formulation the macroscopic current is
evaluated by contracting the propagated density matrix
$\mathbf{P}_\mathbf{k}(t)$ (Section~\ref{sec:propagation}) with
a fixed current operator matrix that is assembled once before the
time loop.
The current operator matrix element in the Kohn-Sham orbital basis
at $\mathbf{k}$-point $\mathbf{k}$ is
\begin{equation}
  \boxed{
  \Jmat{\alpha}_{ab}(\mathbf{k})
  = \underbrace{
      \langle\psi_{a,\mathbf{k}}|
      (-\imath\nabla_\alpha + k_\alpha)
      |\psi_{b,\mathbf{k}}\rangle
    }_{\displaystyle\texttt{VecPHmatrix()}}
  +
  \underbrace{
      \langle\psi_{a,\mathbf{k}}|
      \tfrac{1}{\imath}[\hat{r}_\alpha,\Vnl]
      |\psi_{b,\mathbf{k}}\rangle
    }_{\displaystyle\texttt{CurrentNlpp()}},
  }
  \label{eq:Jmat-full}
\end{equation}
where the code labels refer to \texttt{RMG} routines described in
Section~\ref{sec:implementation}.
The macroscopic $\alpha$-component of the current is then
\begin{equation}
  J_\alpha(t)
  = \sum_{\mathbf{k}} w_\mathbf{k}\,
    \operatorname{Re}\operatorname{Tr}\!\left[
      \mathbf{P}_\mathbf{k}(t)\cdot
      \boldsymbol{\mathcal{J}}^\alpha_\mathbf{k}
    \right],
  \label{eq:J-dm}
\end{equation}
where $w_\mathbf{k}$ are the $\mathbf{k}$-point weights
(with $\sum_\mathbf{k} w_\mathbf{k} = 1$) and the trace runs over
the active orbital subspace.
In the code, eq.~(\ref{eq:J-dm}) is evaluated at each time step by
a single BLAS complex dot product:
\begin{verbatim}
J[alpha] += Re(zdotc(Norb^2, Pn0, 1, Palpha_matrix, 1)) * kweight;
\end{verbatim}
where \texttt{Palpha\_matrix} contains the stored matrix
$\Jmat{\alpha}_{ab}$ (built by \texttt{VecPHmatrix} +
\texttt{CurrentNlpp} before the time loop) and \texttt{Pn0}
contains the current density matrix elements $P_{ab}(t)$.

%--------------------------------------------------------------------
\subsubsection{Hypervirial Relation and Gauge Invariance}
\label{sec:hypervirial}
%--------------------------------------------------------------------

The hypervirial relation is a consistency condition that connects the
velocity (current) operator in the velocity gauge to the dipole
matrix elements in the length gauge.
It is the periodic-system analog of the relation between oscillator
strengths in the length and velocity forms.

In the Heisenberg picture, the Ehrenfest theorem applied to
$\hat{r}_\alpha$ gives:
\begin{equation}
  \langle\psi_a|\hat{v}_\alpha|\psi_b\rangle
  = \langle\psi_a|\tfrac{1}{\imath}[\hat{r}_\alpha,\hat{H}]
    |\psi_b\rangle
  = (\varepsilon_b - \varepsilon_a)
    \langle\psi_a|\hat{r}_\alpha|\psi_b\rangle,
  \label{eq:ehrenfest}
\end{equation}
where $\varepsilon_a$, $\varepsilon_b$ are the ground-state
Kohn-Sham eigenvalues.
For a Hamiltonian with only local potentials, where $\hat{v}_\alpha
= -\imath\nabla_\alpha$, this gives the standard hypervirial relation
\begin{equation}
  \langle\psi_a|(-\imath\nabla_\alpha)|\psi_b\rangle
  = (\varepsilon_b - \varepsilon_a)
    \langle\psi_a|\hat{r}_\alpha|\psi_b\rangle,
  \label{eq:hypervirial-local}
\end{equation}
which states that the velocity gauge matrix element (left side)
equals the length gauge matrix element (right side): both gauges
are equivalent.

When the Hamiltonian contains a nonlocal pseudopotential, the
velocity operator gains the correction
$\frac{1}{\imath}[\hat{r}_\alpha,\Vnl]$
(eq.~\ref{eq:v-full}).
Inserting into eq.~(\ref{eq:ehrenfest}), and rearranging:
\begin{equation}
  \boxed{
  \langle\psi_a|(-\imath\nabla_\alpha)|\psi_b\rangle
  = (\varepsilon_b - \varepsilon_a)
    \langle\psi_a|\hat{r}_\alpha|\psi_b\rangle
  - \langle\psi_a|\tfrac{1}{\imath}
    [\hat{r}_\alpha,\Vnl]|\psi_b\rangle.
  }
  \label{eq:hypervirial-nlpp}
\end{equation}
Equation~(\ref{eq:hypervirial-nlpp}) is the nonlocal-PP-corrected
hypervirial relation, first given by Starace\cite{starace-1971} and
discussed for RT-TDDFT by Mattiat and
Luber.\cite{Mattiat-Luber-2022}
Its physical meaning is that omitting the commutator term from the
velocity operator would give different results in the velocity and
length gauges — that is, wrong spectra.
Including the correction, and only then, restores the gauge
equivalence and ensures that the computed optical conductivity is
independent of the chosen gauge.
Equation~(\ref{eq:hypervirial-nlpp}) also provides a practical
numerical diagnostic: the left-hand side is computed by
\texttt{VecPHmatrix()}, while the right-hand side can be computed
independently from the Kohn-Sham eigenvalues and the dipole matrix;
agreement between the two, up to the grid discretization error,
validates the implementation.

%--------------------------------------------------------------------
\subsubsection{Sign Convention and Factor of $\imath$ in the Code}
\label{sec:imath-convention}
%--------------------------------------------------------------------

We track the factors of $\imath$ explicitly to verify that the code
implements eq.~(\ref{eq:J-dm}) with the correct sign.

The physical current matrix element
$\Jmat{\alpha}_{ab} = \langle\psi_a|(-\imath\nabla_\alpha + k_\alpha
+ \frac{1}{\imath}[\hat{r}_\alpha,\Vnl])|\psi_b\rangle$
is complex-valued in general.
In \texttt{VecPHmatrix()}, the gradient matrix is computed and stored
multiplied by the prefactor \texttt{I\_t} $= \imath$:
\begin{equation}
  \texttt{Pxmatrix}[ab]
  \leftarrow
  \imath\cdot\langle\psi_a|(-\imath\nabla_\alpha)|\psi_b\rangle
  + \imath\cdot k_\alpha\langle\psi_a|\psi_b\rangle
  = \langle\psi_a|(\nabla_\alpha - \imath k_\alpha)|\psi_b\rangle.
  \label{eq:stored-para}
\end{equation}
That is, \texttt{Pxmatrix} stores $\imath\Jmat{\alpha}_{ab}$ from
the paramagnetic part.
\texttt{CurrentNlpp()} adds the nonlocal contribution with the same
prefactor:
\begin{equation}
  \texttt{Pxmatrix}[ab]
  \mathrel{+}=
  \imath\cdot\langle\psi_a|
  \tfrac{1}{\imath}[\hat{r}_\alpha,\Vnl]|\psi_b\rangle
  = \langle\psi_a|[\hat{r}_\alpha,\Vnl]|\psi_b\rangle.
  \label{eq:stored-nlpp}
\end{equation}
Thus \texttt{Pxmatrix} $= \imath\cdot\Jmat{\alpha}_{ab}$ in total.
The macroscopic current is extracted as
\begin{align}
  J_\alpha
  &= \operatorname{Re}\operatorname{Tr}
     \!\left[\mathbf{P}\cdot\texttt{Pxmatrix}\right]
   = \operatorname{Re}\operatorname{Tr}
     \!\left[\mathbf{P}\cdot(\imath\boldsymbol{\mathcal{J}}^\alpha)\right].
  \label{eq:J-from-code}
\end{align}
For a Hermitian density matrix $\mathbf{P} = \mathbf{P}^\dag$
and an anti-Hermitian current operator matrix
$(\boldsymbol{\mathcal{J}}^\alpha)^\dag = -\boldsymbol{\mathcal{J}}^\alpha$
(which follows from the anti-Hermitian nature of $-\imath\nabla$
and of $\frac{1}{\imath}[\hat{r},\Vnl]$%
\footnote{Both $-\imath\nabla$ and $\frac{1}{\imath}[\hat{r},\Vnl]$
are anti-Hermitian operators when the projectors $\beta^{lm}_I$
are real-valued, which is the case for norm-conserving
pseudopotentials in \texttt{RMG}.}),
the trace satisfies
$\operatorname{Tr}[\mathbf{P}\boldsymbol{\mathcal{J}}^\alpha]
= -\overline{\operatorname{Tr}[\mathbf{P}\boldsymbol{\mathcal{J}}^\alpha]}$,
i.e., it is purely imaginary.
Therefore
\begin{equation}
  J_\alpha
  = \operatorname{Re}\!\left[\imath\operatorname{Tr}
    (\mathbf{P}\boldsymbol{\mathcal{J}}^\alpha)\right]
  = -\operatorname{Im}\operatorname{Tr}
    \!\left[\mathbf{P}\boldsymbol{\mathcal{J}}^\alpha\right]
  = \operatorname{Re}\operatorname{Tr}
    \!\left[\mathbf{P}\cdot\imath\boldsymbol{\mathcal{J}}^\alpha\right],
  \label{eq:J-sign-check}
\end{equation}
which is consistent with eq.~(\ref{eq:J-dm}) and confirms that the
sign convention is correct.
The uniform prefactor \texttt{I\_t} applied in both
\texttt{VecPHmatrix()} and \texttt{CurrentNlpp()} ensures that the
paramagnetic and nonlocal contributions enter the current with
the same relative sign, as required by eq.~(\ref{eq:Jmat-full}).



 replaces the placeholder in sec_theory.tex.
%%
%% CONVENTIONS:
%%   Atomic units: hbar = m_e = e = 4*pi*eps_0 = 1
%%   Electron charge = -1  (e > 0 is the elementary charge magnitude)
%%   In atomic units the Hamiltonian is
%%       H = (1/2)(-i*nabla + A)^2 + V_loc + V_nl
%%   Imaginary unit: \imath  (dotless i, consistent with paper 1)
%%   Sign check: NEAT-notes-RMG has a typo (-1/2m kinetic energy);
%%               NEAT-project-notes has the correct +1/2m sign — used here.
%%
%% NOTATION MACROS (all in notation.tex):
%%   \phat   = \hat{\mathbf{p}}  canonical momentum
%%   \pihat  = \hat{\boldsymbol{\pi}}  mechanical momentum
%%   \vhat   = \hat{\mathbf{v}}  velocity operator
%%   \rhat   = \hat{\mathbf{r}}  position
%%   \Vnl    = V^{\mathrm{nl}}
%%   \Vloc   = V^{\mathrm{loc}}
%%   \Avec   = \mathbf{A}
%%   \Jpara  = \mathbf{J}_{\mathrm{para}}
%%   \Jdia   = \mathbf{J}_{\mathrm{dia}}
%%   \Jmat{x}= \mathcal{J}^{x}
%%%%%%%%%%%%%%%%%%%%%%%%%%%%%%%%%%%%%%%%%%%%%%%%%%%%%%%%%%%%%%%%%%%%%%%%%%%%%%%

\subsection{Current Operator and Nonlocal Pseudopotential Contributions}
\label{sec:current}

The central observable in periodic RT-TDDFT is the macroscopic
electronic current density $J_\alpha(t)$, from which the optical
conductivity $\sigma_{\alpha\beta}(\omega)$ and dielectric function
$\varepsilon_{\alpha\beta}(\omega)$ are obtained by Fourier analysis
(Section~\ref{sec:postprocessing}).
For periodic systems the electric dipole moment
$\langle\hat{\mathbf{r}}\rangle$ is ill-defined because the position
operator is incompatible with periodic boundary
conditions,\cite{resta-1998,king-smith-vanderbilt-1993} and the
current is the natural replacement.

We derive the current density operator in four steps: (i) we
distinguish canonical from mechanical momentum and define the velocity
operator via the Heisenberg equation of motion, identifying the role
of nonlocal pseudopotentials; (ii) we rederive the same result from
the continuity equation, decomposing the current into paramagnetic,
diamagnetic, and nonlocal pseudopotential contributions; (iii) we
specialize to Bloch states and derive the $\mathbf{k}$-vector
correction explicitly from the Bloch ansatz; (iv) we express the
macroscopic current in density matrix form and verify the sign
convention in the code.
All equations are in atomic units ($\hbar = m_e = e = 4\pi\varepsilon_0
= 1$) with the electron charge equal to $-1$.

%--------------------------------------------------------------------
\subsubsection{Canonical and Mechanical Momentum}
\label{sec:canon-mech}
%--------------------------------------------------------------------

The Hamiltonian of the Kohn-Sham system in the presence of a
spatially uniform, time-dependent vector potential $\Avec(t)$
(long-wavelength approximation, see Section~\ref{sec:EM-coupling})
is
\begin{equation}
  \hat{H}(t) = \frac{1}{2}
    \!\left(\hat{\mathbf{p}} + \Avec(t)\right)^{\!2}
    + \Vloc(\mathbf{r}) + \Vnl,
  \label{eq:H-vg-2}
\end{equation}
where $\hat{\mathbf{p}} = -\imath\nabla$ is the \emph{canonical}
momentum operator and the sign on $\Avec$ reflects that in atomic
units with electron charge $-1$ the coupling is
$q\Avec = (+1)\cdot(-1)\cdot\Avec = -\Avec$, giving mechanical
momentum $\pihat = \hat{\mathbf{p}} + \Avec$.\footnote{In
Gaussian units the coupling is $q\Avec/c = -e\Avec/c$ for an
electron; in atomic units $c=1/\alpha\approx137$ is absorbed by
setting $e=1$ and the electron charge $-1$, giving
$\hat{\boldsymbol{\pi}} = \hat{\mathbf{p}} + \Avec$ as written.}
The canonical momentum is not gauge invariant and has no direct
physical significance;\cite{Sakurai-2nd_ed} the
\emph{mechanical} (kinetic) momentum
\begin{equation}
  \pihat \equiv \hat{\mathbf{p}} + \Avec(t) = -\imath\nabla + \Avec(t)
  \label{eq:pi-mech}
\end{equation}
is gauge invariant and represents the physical velocity of the
particle.\cite{Sakurai-2nd_ed,starace-1971}
In Gaussian basis set language, this is directly analogous to the
fact that only the kinetic energy integral $\langle\phi_\mu|T|
\phi_\nu\rangle$, not the kinetic canonical momentum itself, enters
the observable transition matrix elements.

Expanding the kinetic term in eq.~(\ref{eq:H-vg-2}):
\begin{align}
  \frac{1}{2}\pihat^2
  &= \frac{1}{2}
     \!\left(-\imath\nabla + \Avec\right)^{\!2}
  \nonumber \\
  &= \underbrace{-\tfrac{1}{2}\nabla^2}_{T_0}
   + \underbrace{\frac{\imath}{2}
     \left(\Avec\cdot\nabla + \nabla\cdot\Avec\right)
     }_{\text{paramagnetic coupling}}
   + \underbrace{\frac{1}{2}\Avec^2}_{\text{diamagnetic}},
  \label{eq:kinetic-expand}
\end{align}
where we use $\nabla\cdot\Avec = 0$ in the Coulomb gauge to
simplify $\Avec\cdot\nabla + \nabla\cdot\Avec = 2\Avec\cdot\nabla$.
In the linear response regime the $\Avec^2$ diamagnetic term is
dropped (it is second order in the applied field).

%--------------------------------------------------------------------
\subsubsection{Velocity Operator from the Heisenberg Equation of Motion}
\label{sec:velocity-eom}
%--------------------------------------------------------------------

The velocity operator is defined via the Heisenberg equation of
motion for the position operator:
\begin{equation}
  \hat{v}_\alpha
  \equiv \frac{1}{\imath}\left[\hat{r}_\alpha,\hat{H}\right].
  \label{eq:heisenberg-v}
\end{equation}
We evaluate eq.~(\ref{eq:heisenberg-v}) for each term in
$\hat{H} = \frac{1}{2}\pihat^2 + \Vloc + \Vnl$ separately.

\paragraph{Kinetic energy contribution.}
We compute $\frac{1}{\imath}[\hat{r}_\alpha, -\frac{1}{2}\nabla^2]$
by acting on an arbitrary test function $g(\mathbf{r})$:
\begin{align}
  \left[\hat{r}_\alpha,\nabla_\alpha^2\right]g
  &= r_\alpha\,\nabla_\alpha^2 g
   - \nabla_\alpha^2(r_\alpha g)
  \nonumber \\
  &= r_\alpha\,\nabla_\alpha^2 g
   - \nabla_\alpha\!\left(g + r_\alpha\nabla_\alpha g\right)
  \nonumber \\
  &= r_\alpha\,\nabla_\alpha^2 g
   - \nabla_\alpha g
   - \nabla_\alpha g
   - r_\alpha\nabla_\alpha^2 g
   = -2\nabla_\alpha g,
  \label{eq:comm-r-nabla2}
\end{align}
where we applied the product rule twice.
The same calculation holds for each Cartesian component, giving
$[\hat{r}_\alpha,\nabla^2] = -2\nabla_\alpha$ (where the right-hand
side involves only the $\alpha$-component of $\nabla$).
Therefore
\begin{equation}
  \frac{1}{\imath}
  \!\left[\hat{r}_\alpha,{-\tfrac{1}{2}\nabla^2}\right]
  = \frac{1}{\imath}\cdot(-\tfrac{1}{2})\cdot(-2\nabla_\alpha)
  = \frac{\nabla_\alpha}{\imath}
  = -\imath\nabla_\alpha = \hat{p}_\alpha.
  \label{eq:comm-r-T}
\end{equation}
Including the full mechanical momentum:
$\frac{1}{\imath}[\hat{r}_\alpha, \frac{1}{2}\pihat^2]
= \hat{p}_\alpha + A_\alpha(t) = \hat{\pi}_\alpha$.

\paragraph{Local potential contribution.}
Because $\Vloc(\mathbf{r})$ is multiplicative (diagonal in position
space),
\begin{equation}
  \left[\hat{r}_\alpha,\Vloc(\mathbf{r})\right]g
  = r_\alpha\,\Vloc g - \Vloc\,r_\alpha g = 0.
  \label{eq:comm-r-Vloc}
\end{equation}
This is the same reason that, in a Gaussian basis, the velocity
matrix elements $\langle\mu|-\imath\nabla|\nu\rangle$ and the
length matrix elements $(\varepsilon_\mu - \varepsilon_\nu)
\langle\mu|\hat{r}|\nu\rangle$ are equal when $V$ is local: both
gauges give the same result.

\paragraph{Nonlocal pseudopotential contribution.}
The nonlocal pseudopotential $\Vnl$ is not diagonal in position
space; its action involves an integral over all space,
\begin{equation}
  \left(\Vnl\psi\right)(\mathbf{r})
  = \int \langle\mathbf{r}|\Vnl|\mathbf{r}'\rangle\,
    \psi(\mathbf{r}')\,\mathrm{d}^3r',
  \label{eq:Vnl-nonlocal-action}
\end{equation}
so $\hat{r}_\alpha$ and $\Vnl$ do not share a common eigenbasis
and their commutator is nonzero.
This is directly analogous to the situation in Gaussian basis sets
with effective core potentials (ECPs), where the ECP commutator
$[\hat{r},V_\mathrm{ECP}]$ also gives a nonzero
contribution.\cite{starace-1971}
The explicit form of $\frac{1}{\imath}[\hat{r}_\alpha,\Vnl]$ for
Kleinman-Bylander separable pseudopotentials is derived in
Appendix~\ref{app:gauge_nlpp}.

\paragraph{Full result.}
Collecting all contributions, the velocity operator is
\begin{equation}
  \boxed{
  \hat{v}_\alpha
  = \frac{1}{\imath}\left[\hat{r}_\alpha,\hat{H}\right]
  = \underbrace{\hat{\pi}_\alpha}_{\displaystyle
    = -\imath\nabla_\alpha + A_\alpha}
  \;+\;
  \underbrace{\frac{1}{\imath}
    \left[\hat{r}_\alpha,\Vnl\right]}_{\text{nonlocal PP correction}}.
  }
  \label{eq:v-full}
\end{equation}
For a purely local Hamiltonian the second term vanishes and
$\hat{v}_\alpha = \hat{\pi}_\alpha$.
The nonlocal correction was first identified by
Starace\cite{starace-1971,starace-1973} in the context of
photoionization from atoms described by model potentials, and is
discussed in detail by Mattiat and Luber\cite{Mattiat-Luber-2022}
for RT-TDDFT with nonlocal pseudopotentials.
Its physical origin is that the electron ``feels'' the nonlocal
potential at points $\mathbf{r}'$ away from its instantaneous
position $\mathbf{r}$, so its velocity cannot be expressed through
the local gradient alone.

%--------------------------------------------------------------------
\subsubsection{Current Density from the Continuity Equation}
\label{sec:current-continuity}
%--------------------------------------------------------------------

We now derive the same current density by an independent route
through the continuity equation, which provides a consistency check
and clarifies the physical role of each contribution.

The probability density is $\rho(\mathbf{r},t) = \psi^*\psi$.
Differentiating with respect to time and substituting the TDKS
equation $\imath\partial_t\psi = \hat{H}\psi$ and its Hermitian
conjugate $-\imath\partial_t\psi^* = \hat{H}\psi^*$:
\begin{align}
  \frac{\partial\rho}{\partial t}
  &= \dot\psi^*\psi + \psi^*\dot\psi
   = \imath\!\left[(\hat{H}\psi^*)\psi
                  - \psi^*(\hat{H}\psi)\right].
  \label{eq:drho-H}
\end{align}
We evaluate the right-hand side term by term.

\paragraph{Local potential: vanishes.}
$\imath[(\Vloc\psi^*)\psi - \psi^*(\Vloc\psi)]
= \imath\Vloc(\psi^*\psi - \psi^*\psi) = 0$.

\paragraph{Kinetic energy: paramagnetic and diamagnetic currents.}
Using the Coulomb gauge $\nabla\cdot\Avec = 0$ and abbreviating
$\hat{D} \equiv -\imath\nabla + \Avec$:
\begin{align}
  \imath\!\left[(\hat{T}\psi^*)\psi - \psi^*(\hat{T}\psi)\right]
  &= \frac{\imath}{2}\!\left[
     \overline{(\hat{D}^2\psi)}\cdot\psi
     - \psi^*\cdot(\hat{D}^2\psi)
     \right]
  \nonumber \\
  &= -\nabla\cdot\!\underbrace{\left[
     \frac{1}{2}\operatorname{Re}\!\left(
     \psi^*\hat{D}\psi\right)\right]}_{\displaystyle
     = \Jpara + \Jdia},
  \label{eq:kin-continuity}
\end{align}
where the last equality follows from integration by parts applied
to $\nabla\cdot(\psi^*\hat{D}\psi)$.
Expanding $\hat{D} = -\imath\nabla + \Avec$ inside the
$\operatorname{Re}[\cdot]$:
\begin{align}
  \Jpara(\mathbf{r},t)
  &= \operatorname{Re}\!\left[\psi^*(-\imath\nabla)\psi\right],
  \label{eq:Jpara-def}
  \\
  \Jdia(\mathbf{r},t)
  &= \Avec(t)\,|\psi(\mathbf{r},t)|^2.
  \label{eq:Jdia-def}
\end{align}

\paragraph{Nonlocal pseudopotential: nonlocal current.}
The nonlocal term contributes
\begin{equation}
  \imath\!\left[(\Vnl\psi^*)\psi - \psi^*(\Vnl\psi)\right]
  = -\nabla\cdot\mathbf{J}_\mathrm{nl}(\mathbf{r},t),
  \label{eq:Jnl-continuity}
\end{equation}
where $\mathbf{J}_\mathrm{nl}$ is the nonlocal current density
whose explicit form in terms of the Kleinman-Bylander projectors
is derived in Appendix~\ref{app:Jnl-expression}.

\paragraph{Total current density.}
Comparing with the continuity equation
$\partial_t\rho + \nabla\cdot\mathbf{J} = 0$ and collecting
all contributions:
\begin{equation}
  \mathbf{J}(\mathbf{r},t)
  = \underbrace{
      \operatorname{Re}\!\left[\psi^*(-\imath\nabla)\psi\right]
    }_{\Jpara}
  + \underbrace{
      \Avec(t)|\psi|^2
    }_{\Jdia}
  + \underbrace{
      \operatorname{Re}\!\left[
        \psi^*\tfrac{1}{\imath}[\hat{\mathbf{r}},\Vnl]\psi
      \right]
    }_{\mathbf{J}_\mathrm{nl}}.
  \label{eq:J-total}
\end{equation}
This is fully consistent with eq.~(\ref{eq:v-full}): the three
contributions map one-to-one onto the mechanical momentum
$\hat{\pi}_\alpha$ (split into its $\hat{p}_\alpha$ and $A_\alpha$
parts) plus the nonlocal correction.

%--------------------------------------------------------------------
\subsubsection{Paramagnetic and Diamagnetic Contributions:
  Physical Roles}
\label{sec:para-dia}
%--------------------------------------------------------------------

The \emph{paramagnetic} current $\Jpara$ is independent of the
external field $\Avec$ and represents the intrinsic probability
current of the time-evolving quantum state.
It oscillates in response to the delta-kick perturbation and, after
Fourier transformation, yields the optical conductivity.

The \emph{diamagnetic} current $\Jdia = \Avec(t)|\psi|^2$ is
proportional to the applied vector potential.
In the perturbative kick scheme used in \texttt{RMG},%
\cite{yabana-bertsch-1996,jakowski-rmg-tddft-2025}
the vector potential is a delta-function impulse:
$\Avec(t) = \Avec_0\,\delta(t)$, with $\Avec(t) = 0$ for all
$t > 0$.
Consequently:
\begin{equation}
  \Jdia(\mathbf{r},t>0) = \Avec(t>0)\,|\psi|^2 = 0.
  \label{eq:Jdia-zero}
\end{equation}
The diamagnetic current vanishes identically during the simulation
after the kick.
This is not an approximation: within linear response theory the
complete first-order optical response is encoded in the oscillations
of $\Jpara(t>0)$ and $\mathbf{J}_\mathrm{nl}(t>0)$,%
\cite{yabana-bertsch-1996}
so the macroscopic current that is recorded and Fourier-transformed
is
\begin{equation}
  J_\alpha(t>0)
  = \int_\Omega\!
    \left[\Jpara(\mathbf{r},t)
    + \mathbf{J}_{\mathrm{nl}}(\mathbf{r},t)\right]_\alpha
    \mathrm{d}^3r,
  \label{eq:J-sim}
\end{equation}
where the integral is over the unit cell $\Omega$.

%--------------------------------------------------------------------
\subsubsection{Current Operator for Bloch States:
  the $\mathbf{k}$-Vector Correction}
\label{sec:current-bloch}
%--------------------------------------------------------------------

For a periodic crystal, the Kohn-Sham orbitals satisfy Bloch's
theorem.
We derive how the current operator is modified when expressed in
terms of the cell-periodic part $u_{n,\mathbf{k}}$.

A Bloch state at band index $n$ and crystal momentum $\mathbf{k}$
is written as
\begin{equation}
  \psi_{n,\mathbf{k}}(\mathbf{r},t)
  = e^{\imath\mathbf{k}\cdot\mathbf{r}}\,
    u_{n,\mathbf{k}}(\mathbf{r},t),
  \label{eq:bloch-def}
\end{equation}
where $u_{n,\mathbf{k}}(\mathbf{r}+\mathbf{R},t)
= u_{n,\mathbf{k}}(\mathbf{r},t)$ for all lattice vectors
$\mathbf{R}$.
Applying the canonical momentum operator $-\imath\nabla_\alpha$ to
the full Bloch state and using the product rule:
\begin{align}
  (-\imath\nabla_\alpha)\,\psi_{n,\mathbf{k}}
  &= (-\imath\nabla_\alpha)\!\left[
     e^{\imath\mathbf{k}\cdot\mathbf{r}}\,u_{n,\mathbf{k}}
     \right]
  \nonumber \\
  &= (-\imath)\!\left[
     \imath k_\alpha\,
     e^{\imath\mathbf{k}\cdot\mathbf{r}}\,u_{n,\mathbf{k}}
     + e^{\imath\mathbf{k}\cdot\mathbf{r}}\,
       \nabla_\alpha u_{n,\mathbf{k}}
     \right]
  \nonumber \\
  &= e^{\imath\mathbf{k}\cdot\mathbf{r}}
     \!\left[k_\alpha + (-\imath\nabla_\alpha)\right]
     u_{n,\mathbf{k}}.
  \label{eq:nabla-bloch}
\end{align}
The Bloch phase factor $e^{\imath\mathbf{k}\cdot\mathbf{r}}$
cancels between bra and ket in a matrix element over the unit cell,
so the matrix element of the canonical momentum between two Bloch
states at the same $\mathbf{k}$-point is
\begin{align}
  p^\alpha_{ab}(\mathbf{k})
  &\equiv \langle\psi_{a,\mathbf{k}}|(-\imath\nabla_\alpha)|
    \psi_{b,\mathbf{k}}\rangle
  \nonumber \\
  &= \int_\Omega\!
     e^{-\imath\mathbf{k}\cdot\mathbf{r}}\,u_{a,\mathbf{k}}^*\cdot
     e^{\imath\mathbf{k}\cdot\mathbf{r}}
     \!\left(-\imath\nabla_\alpha + k_\alpha\right)
     u_{b,\mathbf{k}}\,\mathrm{d}^3r
  \nonumber \\
  &= \langle u_{a,\mathbf{k}}|
     \!\left(-\imath\nabla_\alpha + k_\alpha\right)
     |u_{b,\mathbf{k}}\rangle_\Omega.
  \label{eq:para-bloch-me}
\end{align}
The $k_\alpha$ shift arises purely from the Bloch factor and is
absent for $\Gamma$-point ($\mathbf{k} = \mathbf{0}$) calculations.
For finite $\mathbf{k}$-point sampling it must be included to obtain
the correct optical response; its omission would give wrong
intensities proportional to $|\mathbf{k}|^2$ at each $\mathbf{k}$-point.

The full current density summed over all occupied Bloch states is
\begin{equation}
  \mathbf{J}(\mathbf{r},t)
  = \frac{1}{N_k}\!\sum_{n,\mathbf{k}}^{\mathrm{occ}}
    f_{n,\mathbf{k}}\,
    \operatorname{Re}\!\left[
      u_{n,\mathbf{k}}^*\!\left(
      -\imath\nabla + \mathbf{k}
      + \frac{1}{\imath}[\hat{\mathbf{r}},\Vnl]
      \right)u_{n,\mathbf{k}}
    \right],
  \label{eq:J-bloch-total}
\end{equation}
where $N_k$ is the number of $\mathbf{k}$-points and $f_{n,\mathbf{k}}$
are the occupation numbers.
Equation~(\ref{eq:J-bloch-total}) reproduces eq.~(8) of Yabana and
Bertsch\cite{yabana-bertsch-1996} and eq.~(26) of Mattiat and
Luber,\cite{Mattiat-Luber-2022} specialized to the Coulomb gauge
and the velocity gauge with a spatially uniform $\Avec(t)$.

%--------------------------------------------------------------------
\subsubsection{Current Operator Matrix and Density Matrix Formulation}
\label{sec:current-dm-sec}
%--------------------------------------------------------------------

In the density matrix formulation the macroscopic current is
evaluated by contracting the propagated density matrix
$\mathbf{P}_\mathbf{k}(t)$ (Section~\ref{sec:propagation}) with
a fixed current operator matrix that is assembled once before the
time loop.
The current operator matrix element in the Kohn-Sham orbital basis
at $\mathbf{k}$-point $\mathbf{k}$ is
\begin{equation}
  \boxed{
  \Jmat{\alpha}_{ab}(\mathbf{k})
  = \underbrace{
      \langle\psi_{a,\mathbf{k}}|
      (-\imath\nabla_\alpha + k_\alpha)
      |\psi_{b,\mathbf{k}}\rangle
    }_{\displaystyle\texttt{VecPHmatrix()}}
  +
  \underbrace{
      \langle\psi_{a,\mathbf{k}}|
      \tfrac{1}{\imath}[\hat{r}_\alpha,\Vnl]
      |\psi_{b,\mathbf{k}}\rangle
    }_{\displaystyle\texttt{CurrentNlpp()}},
  }
  \label{eq:Jmat-full}
\end{equation}
where the code labels refer to \texttt{RMG} routines described in
Section~\ref{sec:implementation}.
The macroscopic $\alpha$-component of the current is then
\begin{equation}
  J_\alpha(t)
  = \sum_{\mathbf{k}} w_\mathbf{k}\,
    \operatorname{Re}\operatorname{Tr}\!\left[
      \mathbf{P}_\mathbf{k}(t)\cdot
      \boldsymbol{\mathcal{J}}^\alpha_\mathbf{k}
    \right],
  \label{eq:J-dm}
\end{equation}
where $w_\mathbf{k}$ are the $\mathbf{k}$-point weights
(with $\sum_\mathbf{k} w_\mathbf{k} = 1$) and the trace runs over
the active orbital subspace.
In the code, eq.~(\ref{eq:J-dm}) is evaluated at each time step by
a single BLAS complex dot product:
\begin{verbatim}
J[alpha] += Re(zdotc(Norb^2, Pn0, 1, Palpha_matrix, 1)) * kweight;
\end{verbatim}
where \texttt{Palpha\_matrix} contains the stored matrix
$\Jmat{\alpha}_{ab}$ (built by \texttt{VecPHmatrix} +
\texttt{CurrentNlpp} before the time loop) and \texttt{Pn0}
contains the current density matrix elements $P_{ab}(t)$.

%--------------------------------------------------------------------
\subsubsection{Hypervirial Relation and Gauge Invariance}
\label{sec:hypervirial}
%--------------------------------------------------------------------

The hypervirial relation is a consistency condition that connects the
velocity (current) operator in the velocity gauge to the dipole
matrix elements in the length gauge.
It is the periodic-system analog of the relation between oscillator
strengths in the length and velocity forms.

In the Heisenberg picture, the Ehrenfest theorem applied to
$\hat{r}_\alpha$ gives:
\begin{equation}
  \langle\psi_a|\hat{v}_\alpha|\psi_b\rangle
  = \langle\psi_a|\tfrac{1}{\imath}[\hat{r}_\alpha,\hat{H}]
    |\psi_b\rangle
  = (\varepsilon_b - \varepsilon_a)
    \langle\psi_a|\hat{r}_\alpha|\psi_b\rangle,
  \label{eq:ehrenfest}
\end{equation}
where $\varepsilon_a$, $\varepsilon_b$ are the ground-state
Kohn-Sham eigenvalues.
For a Hamiltonian with only local potentials, where $\hat{v}_\alpha
= -\imath\nabla_\alpha$, this gives the standard hypervirial relation
\begin{equation}
  \langle\psi_a|(-\imath\nabla_\alpha)|\psi_b\rangle
  = (\varepsilon_b - \varepsilon_a)
    \langle\psi_a|\hat{r}_\alpha|\psi_b\rangle,
  \label{eq:hypervirial-local}
\end{equation}
which states that the velocity gauge matrix element (left side)
equals the length gauge matrix element (right side): both gauges
are equivalent.

When the Hamiltonian contains a nonlocal pseudopotential, the
velocity operator gains the correction
$\frac{1}{\imath}[\hat{r}_\alpha,\Vnl]$
(eq.~\ref{eq:v-full}).
Inserting into eq.~(\ref{eq:ehrenfest}), and rearranging:
\begin{equation}
  \boxed{
  \langle\psi_a|(-\imath\nabla_\alpha)|\psi_b\rangle
  = (\varepsilon_b - \varepsilon_a)
    \langle\psi_a|\hat{r}_\alpha|\psi_b\rangle
  - \langle\psi_a|\tfrac{1}{\imath}
    [\hat{r}_\alpha,\Vnl]|\psi_b\rangle.
  }
  \label{eq:hypervirial-nlpp}
\end{equation}
Equation~(\ref{eq:hypervirial-nlpp}) is the nonlocal-PP-corrected
hypervirial relation, first given by Starace\cite{starace-1971} and
discussed for RT-TDDFT by Mattiat and
Luber.\cite{Mattiat-Luber-2022}
Its physical meaning is that omitting the commutator term from the
velocity operator would give different results in the velocity and
length gauges — that is, wrong spectra.
Including the correction, and only then, restores the gauge
equivalence and ensures that the computed optical conductivity is
independent of the chosen gauge.
Equation~(\ref{eq:hypervirial-nlpp}) also provides a practical
numerical diagnostic: the left-hand side is computed by
\texttt{VecPHmatrix()}, while the right-hand side can be computed
independently from the Kohn-Sham eigenvalues and the dipole matrix;
agreement between the two, up to the grid discretization error,
validates the implementation.

%--------------------------------------------------------------------
\subsubsection{Sign Convention and Factor of $\imath$ in the Code}
\label{sec:imath-convention}
%--------------------------------------------------------------------

We track the factors of $\imath$ explicitly to verify that the code
implements eq.~(\ref{eq:J-dm}) with the correct sign.

The physical current matrix element
$\Jmat{\alpha}_{ab} = \langle\psi_a|(-\imath\nabla_\alpha + k_\alpha
+ \frac{1}{\imath}[\hat{r}_\alpha,\Vnl])|\psi_b\rangle$
is complex-valued in general.
In \texttt{VecPHmatrix()}, the gradient matrix is computed and stored
multiplied by the prefactor \texttt{I\_t} $= \imath$:
\begin{equation}
  \texttt{Pxmatrix}[ab]
  \leftarrow
  \imath\cdot\langle\psi_a|(-\imath\nabla_\alpha)|\psi_b\rangle
  + \imath\cdot k_\alpha\langle\psi_a|\psi_b\rangle
  = \langle\psi_a|(\nabla_\alpha - \imath k_\alpha)|\psi_b\rangle.
  \label{eq:stored-para}
\end{equation}
That is, \texttt{Pxmatrix} stores $\imath\Jmat{\alpha}_{ab}$ from
the paramagnetic part.
\texttt{CurrentNlpp()} adds the nonlocal contribution with the same
prefactor:
\begin{equation}
  \texttt{Pxmatrix}[ab]
  \mathrel{+}=
  \imath\cdot\langle\psi_a|
  \tfrac{1}{\imath}[\hat{r}_\alpha,\Vnl]|\psi_b\rangle
  = \langle\psi_a|[\hat{r}_\alpha,\Vnl]|\psi_b\rangle.
  \label{eq:stored-nlpp}
\end{equation}
Thus \texttt{Pxmatrix} $= \imath\cdot\Jmat{\alpha}_{ab}$ in total.
The macroscopic current is extracted as
\begin{align}
  J_\alpha
  &= \operatorname{Re}\operatorname{Tr}
     \!\left[\mathbf{P}\cdot\texttt{Pxmatrix}\right]
   = \operatorname{Re}\operatorname{Tr}
     \!\left[\mathbf{P}\cdot(\imath\boldsymbol{\mathcal{J}}^\alpha)\right].
  \label{eq:J-from-code}
\end{align}
For a Hermitian density matrix $\mathbf{P} = \mathbf{P}^\dag$
and an anti-Hermitian current operator matrix
$(\boldsymbol{\mathcal{J}}^\alpha)^\dag = -\boldsymbol{\mathcal{J}}^\alpha$
(which follows from the anti-Hermitian nature of $-\imath\nabla$
and of $\frac{1}{\imath}[\hat{r},\Vnl]$%
\footnote{Both $-\imath\nabla$ and $\frac{1}{\imath}[\hat{r},\Vnl]$
are anti-Hermitian operators when the projectors $\beta^{lm}_I$
are real-valued, which is the case for norm-conserving
pseudopotentials in \texttt{RMG}.}),
the trace satisfies
$\operatorname{Tr}[\mathbf{P}\boldsymbol{\mathcal{J}}^\alpha]
= -\overline{\operatorname{Tr}[\mathbf{P}\boldsymbol{\mathcal{J}}^\alpha]}$,
i.e., it is purely imaginary.
Therefore
\begin{equation}
  J_\alpha
  = \operatorname{Re}\!\left[\imath\operatorname{Tr}
    (\mathbf{P}\boldsymbol{\mathcal{J}}^\alpha)\right]
  = -\operatorname{Im}\operatorname{Tr}
    \!\left[\mathbf{P}\boldsymbol{\mathcal{J}}^\alpha\right]
  = \operatorname{Re}\operatorname{Tr}
    \!\left[\mathbf{P}\cdot\imath\boldsymbol{\mathcal{J}}^\alpha\right],
  \label{eq:J-sign-check}
\end{equation}
which is consistent with eq.~(\ref{eq:J-dm}) and confirms that the
sign convention is correct.
The uniform prefactor \texttt{I\_t} applied in both
\texttt{VecPHmatrix()} and \texttt{CurrentNlpp()} ensures that the
paramagnetic and nonlocal contributions enter the current with
the same relative sign, as required by eq.~(\ref{eq:Jmat-full}).



 replaces the placeholder in sec_theory.tex.
%%
%% CONVENTIONS:
%%   Atomic units: hbar = m_e = e = 4*pi*eps_0 = 1
%%   Electron charge = -1  (e > 0 is the elementary charge magnitude)
%%   In atomic units the Hamiltonian is
%%       H = (1/2)(-i*nabla + A)^2 + V_loc + V_nl
%%   Imaginary unit: \imath  (dotless i, consistent with paper 1)
%%   Sign check: NEAT-notes-RMG has a typo (-1/2m kinetic energy);
%%               NEAT-project-notes has the correct +1/2m sign — used here.
%%
%% NOTATION MACROS (all in notation.tex):
%%   \phat   = \hat{\mathbf{p}}  canonical momentum
%%   \pihat  = \hat{\boldsymbol{\pi}}  mechanical momentum
%%   \vhat   = \hat{\mathbf{v}}  velocity operator
%%   \rhat   = \hat{\mathbf{r}}  position
%%   \Vnl    = V^{\mathrm{nl}}
%%   \Vloc   = V^{\mathrm{loc}}
%%   \Avec   = \mathbf{A}
%%   \Jpara  = \mathbf{J}_{\mathrm{para}}
%%   \Jdia   = \mathbf{J}_{\mathrm{dia}}
%%   \Jmat{x}= \mathcal{J}^{x}
%%%%%%%%%%%%%%%%%%%%%%%%%%%%%%%%%%%%%%%%%%%%%%%%%%%%%%%%%%%%%%%%%%%%%%%%%%%%%%%

\subsection{Current Operator and Nonlocal Pseudopotential Contributions}
\label{sec:current}

The central observable in periodic RT-TDDFT is the macroscopic
electronic current density $J_\alpha(t)$, from which the optical
conductivity $\sigma_{\alpha\beta}(\omega)$ and dielectric function
$\varepsilon_{\alpha\beta}(\omega)$ are obtained by Fourier analysis
(Section~\ref{sec:postprocessing}).
For periodic systems the electric dipole moment
$\langle\hat{\mathbf{r}}\rangle$ is ill-defined because the position
operator is incompatible with periodic boundary
conditions,\cite{resta-1998,king-smith-vanderbilt-1993} and the
current is the natural replacement.

We derive the current density operator in four steps: (i) we
distinguish canonical from mechanical momentum and define the velocity
operator via the Heisenberg equation of motion, identifying the role
of nonlocal pseudopotentials; (ii) we rederive the same result from
the continuity equation, decomposing the current into paramagnetic,
diamagnetic, and nonlocal pseudopotential contributions; (iii) we
specialize to Bloch states and derive the $\mathbf{k}$-vector
correction explicitly from the Bloch ansatz; (iv) we express the
macroscopic current in density matrix form and verify the sign
convention in the code.
All equations are in atomic units ($\hbar = m_e = e = 4\pi\varepsilon_0
= 1$) with the electron charge equal to $-1$.

%--------------------------------------------------------------------
\subsubsection{Canonical and Mechanical Momentum}
\label{sec:canon-mech}
%--------------------------------------------------------------------

The Hamiltonian of the Kohn-Sham system in the presence of a
spatially uniform, time-dependent vector potential $\Avec(t)$
(long-wavelength approximation, see Section~\ref{sec:EM-coupling})
is
\begin{equation}
  \hat{H}(t) = \frac{1}{2}
    \!\left(\hat{\mathbf{p}} + \Avec(t)\right)^{\!2}
    + \Vloc(\mathbf{r}) + \Vnl,
  \label{eq:H-vg-2}
\end{equation}
where $\hat{\mathbf{p}} = -\imath\nabla$ is the \emph{canonical}
momentum operator and the sign on $\Avec$ reflects that in atomic
units with electron charge $-1$ the coupling is
$q\Avec = (+1)\cdot(-1)\cdot\Avec = -\Avec$, giving mechanical
momentum $\pihat = \hat{\mathbf{p}} + \Avec$.\footnote{In
Gaussian units the coupling is $q\Avec/c = -e\Avec/c$ for an
electron; in atomic units $c=1/\alpha\approx137$ is absorbed by
setting $e=1$ and the electron charge $-1$, giving
$\hat{\boldsymbol{\pi}} = \hat{\mathbf{p}} + \Avec$ as written.}
The canonical momentum is not gauge invariant and has no direct
physical significance;\cite{Sakurai-2nd_ed} the
\emph{mechanical} (kinetic) momentum
\begin{equation}
  \pihat \equiv \hat{\mathbf{p}} + \Avec(t) = -\imath\nabla + \Avec(t)
  \label{eq:pi-mech}
\end{equation}
is gauge invariant and represents the physical velocity of the
particle.\cite{Sakurai-2nd_ed,starace-1971}
In Gaussian basis set language, this is directly analogous to the
fact that only the kinetic energy integral $\langle\phi_\mu|T|
\phi_\nu\rangle$, not the kinetic canonical momentum itself, enters
the observable transition matrix elements.

Expanding the kinetic term in eq.~(\ref{eq:H-vg-2}):
\begin{align}
  \frac{1}{2}\pihat^2
  &= \frac{1}{2}
     \!\left(-\imath\nabla + \Avec\right)^{\!2}
  \nonumber \\
  &= \underbrace{-\tfrac{1}{2}\nabla^2}_{T_0}
   + \underbrace{\frac{\imath}{2}
     \left(\Avec\cdot\nabla + \nabla\cdot\Avec\right)
     }_{\text{paramagnetic coupling}}
   + \underbrace{\frac{1}{2}\Avec^2}_{\text{diamagnetic}},
  \label{eq:kinetic-expand}
\end{align}
where we use $\nabla\cdot\Avec = 0$ in the Coulomb gauge to
simplify $\Avec\cdot\nabla + \nabla\cdot\Avec = 2\Avec\cdot\nabla$.
In the linear response regime the $\Avec^2$ diamagnetic term is
dropped (it is second order in the applied field).

%--------------------------------------------------------------------
\subsubsection{Velocity Operator from the Heisenberg Equation of Motion}
\label{sec:velocity-eom}
%--------------------------------------------------------------------

The velocity operator is defined via the Heisenberg equation of
motion for the position operator:
\begin{equation}
  \hat{v}_\alpha
  \equiv \frac{1}{\imath}\left[\hat{r}_\alpha,\hat{H}\right].
  \label{eq:heisenberg-v}
\end{equation}
We evaluate eq.~(\ref{eq:heisenberg-v}) for each term in
$\hat{H} = \frac{1}{2}\pihat^2 + \Vloc + \Vnl$ separately.

\paragraph{Kinetic energy contribution.}
We compute $\frac{1}{\imath}[\hat{r}_\alpha, -\frac{1}{2}\nabla^2]$
by acting on an arbitrary test function $g(\mathbf{r})$:
\begin{align}
  \left[\hat{r}_\alpha,\nabla_\alpha^2\right]g
  &= r_\alpha\,\nabla_\alpha^2 g
   - \nabla_\alpha^2(r_\alpha g)
  \nonumber \\
  &= r_\alpha\,\nabla_\alpha^2 g
   - \nabla_\alpha\!\left(g + r_\alpha\nabla_\alpha g\right)
  \nonumber \\
  &= r_\alpha\,\nabla_\alpha^2 g
   - \nabla_\alpha g
   - \nabla_\alpha g
   - r_\alpha\nabla_\alpha^2 g
   = -2\nabla_\alpha g,
  \label{eq:comm-r-nabla2}
\end{align}
where we applied the product rule twice.
The same calculation holds for each Cartesian component, giving
$[\hat{r}_\alpha,\nabla^2] = -2\nabla_\alpha$ (where the right-hand
side involves only the $\alpha$-component of $\nabla$).
Therefore
\begin{equation}
  \frac{1}{\imath}
  \!\left[\hat{r}_\alpha,{-\tfrac{1}{2}\nabla^2}\right]
  = \frac{1}{\imath}\cdot(-\tfrac{1}{2})\cdot(-2\nabla_\alpha)
  = \frac{\nabla_\alpha}{\imath}
  = -\imath\nabla_\alpha = \hat{p}_\alpha.
  \label{eq:comm-r-T}
\end{equation}
Including the full mechanical momentum:
$\frac{1}{\imath}[\hat{r}_\alpha, \frac{1}{2}\pihat^2]
= \hat{p}_\alpha + A_\alpha(t) = \hat{\pi}_\alpha$.

\paragraph{Local potential contribution.}
Because $\Vloc(\mathbf{r})$ is multiplicative (diagonal in position
space),
\begin{equation}
  \left[\hat{r}_\alpha,\Vloc(\mathbf{r})\right]g
  = r_\alpha\,\Vloc g - \Vloc\,r_\alpha g = 0.
  \label{eq:comm-r-Vloc}
\end{equation}
This is the same reason that, in a Gaussian basis, the velocity
matrix elements $\langle\mu|-\imath\nabla|\nu\rangle$ and the
length matrix elements $(\varepsilon_\mu - \varepsilon_\nu)
\langle\mu|\hat{r}|\nu\rangle$ are equal when $V$ is local: both
gauges give the same result.

\paragraph{Nonlocal pseudopotential contribution.}
The nonlocal pseudopotential $\Vnl$ is not diagonal in position
space; its action involves an integral over all space,
\begin{equation}
  \left(\Vnl\psi\right)(\mathbf{r})
  = \int \langle\mathbf{r}|\Vnl|\mathbf{r}'\rangle\,
    \psi(\mathbf{r}')\,\mathrm{d}^3r',
  \label{eq:Vnl-nonlocal-action}
\end{equation}
so $\hat{r}_\alpha$ and $\Vnl$ do not share a common eigenbasis
and their commutator is nonzero.
This is directly analogous to the situation in Gaussian basis sets
with effective core potentials (ECPs), where the ECP commutator
$[\hat{r},V_\mathrm{ECP}]$ also gives a nonzero
contribution.\cite{starace-1971}
The explicit form of $\frac{1}{\imath}[\hat{r}_\alpha,\Vnl]$ for
Kleinman-Bylander separable pseudopotentials is derived in
Appendix~\ref{app:gauge_nlpp}.

\paragraph{Full result.}
Collecting all contributions, the velocity operator is
\begin{equation}
  \boxed{
  \hat{v}_\alpha
  = \frac{1}{\imath}\left[\hat{r}_\alpha,\hat{H}\right]
  = \underbrace{\hat{\pi}_\alpha}_{\displaystyle
    = -\imath\nabla_\alpha + A_\alpha}
  \;+\;
  \underbrace{\frac{1}{\imath}
    \left[\hat{r}_\alpha,\Vnl\right]}_{\text{nonlocal PP correction}}.
  }
  \label{eq:v-full}
\end{equation}
For a purely local Hamiltonian the second term vanishes and
$\hat{v}_\alpha = \hat{\pi}_\alpha$.
The nonlocal correction was first identified by
Starace\cite{starace-1971,starace-1973} in the context of
photoionization from atoms described by model potentials, and is
discussed in detail by Mattiat and Luber\cite{Mattiat-Luber-2022}
for RT-TDDFT with nonlocal pseudopotentials.
Its physical origin is that the electron ``feels'' the nonlocal
potential at points $\mathbf{r}'$ away from its instantaneous
position $\mathbf{r}$, so its velocity cannot be expressed through
the local gradient alone.

%--------------------------------------------------------------------
\subsubsection{Current Density from the Continuity Equation}
\label{sec:current-continuity}
%--------------------------------------------------------------------

We now derive the same current density by an independent route
through the continuity equation, which provides a consistency check
and clarifies the physical role of each contribution.

The probability density is $\rho(\mathbf{r},t) = \psi^*\psi$.
Differentiating with respect to time and substituting the TDKS
equation $\imath\partial_t\psi = \hat{H}\psi$ and its Hermitian
conjugate $-\imath\partial_t\psi^* = \hat{H}\psi^*$:
\begin{align}
  \frac{\partial\rho}{\partial t}
  &= \dot\psi^*\psi + \psi^*\dot\psi
   = \imath\!\left[(\hat{H}\psi^*)\psi
                  - \psi^*(\hat{H}\psi)\right].
  \label{eq:drho-H}
\end{align}
We evaluate the right-hand side term by term.

\paragraph{Local potential: vanishes.}
$\imath[(\Vloc\psi^*)\psi - \psi^*(\Vloc\psi)]
= \imath\Vloc(\psi^*\psi - \psi^*\psi) = 0$.

\paragraph{Kinetic energy: paramagnetic and diamagnetic currents.}
Using the Coulomb gauge $\nabla\cdot\Avec = 0$ and abbreviating
$\hat{D} \equiv -\imath\nabla + \Avec$:
\begin{align}
  \imath\!\left[(\hat{T}\psi^*)\psi - \psi^*(\hat{T}\psi)\right]
  &= \frac{\imath}{2}\!\left[
     \overline{(\hat{D}^2\psi)}\cdot\psi
     - \psi^*\cdot(\hat{D}^2\psi)
     \right]
  \nonumber \\
  &= -\nabla\cdot\!\underbrace{\left[
     \frac{1}{2}\operatorname{Re}\!\left(
     \psi^*\hat{D}\psi\right)\right]}_{\displaystyle
     = \Jpara + \Jdia},
  \label{eq:kin-continuity}
\end{align}
where the last equality follows from integration by parts applied
to $\nabla\cdot(\psi^*\hat{D}\psi)$.
Expanding $\hat{D} = -\imath\nabla + \Avec$ inside the
$\operatorname{Re}[\cdot]$:
\begin{align}
  \Jpara(\mathbf{r},t)
  &= \operatorname{Re}\!\left[\psi^*(-\imath\nabla)\psi\right],
  \label{eq:Jpara-def}
  \\
  \Jdia(\mathbf{r},t)
  &= \Avec(t)\,|\psi(\mathbf{r},t)|^2.
  \label{eq:Jdia-def}
\end{align}

\paragraph{Nonlocal pseudopotential: nonlocal current.}
The nonlocal term contributes
\begin{equation}
  \imath\!\left[(\Vnl\psi^*)\psi - \psi^*(\Vnl\psi)\right]
  = -\nabla\cdot\mathbf{J}_\mathrm{nl}(\mathbf{r},t),
  \label{eq:Jnl-continuity}
\end{equation}
where $\mathbf{J}_\mathrm{nl}$ is the nonlocal current density
whose explicit form in terms of the Kleinman-Bylander projectors
is derived in Appendix~\ref{app:Jnl-expression}.

\paragraph{Total current density.}
Comparing with the continuity equation
$\partial_t\rho + \nabla\cdot\mathbf{J} = 0$ and collecting
all contributions:
\begin{equation}
  \mathbf{J}(\mathbf{r},t)
  = \underbrace{
      \operatorname{Re}\!\left[\psi^*(-\imath\nabla)\psi\right]
    }_{\Jpara}
  + \underbrace{
      \Avec(t)|\psi|^2
    }_{\Jdia}
  + \underbrace{
      \operatorname{Re}\!\left[
        \psi^*\tfrac{1}{\imath}[\hat{\mathbf{r}},\Vnl]\psi
      \right]
    }_{\mathbf{J}_\mathrm{nl}}.
  \label{eq:J-total}
\end{equation}
This is fully consistent with eq.~(\ref{eq:v-full}): the three
contributions map one-to-one onto the mechanical momentum
$\hat{\pi}_\alpha$ (split into its $\hat{p}_\alpha$ and $A_\alpha$
parts) plus the nonlocal correction.

%--------------------------------------------------------------------
\subsubsection{Paramagnetic and Diamagnetic Contributions:
  Physical Roles}
\label{sec:para-dia}
%--------------------------------------------------------------------

The \emph{paramagnetic} current $\Jpara$ is independent of the
external field $\Avec$ and represents the intrinsic probability
current of the time-evolving quantum state.
It oscillates in response to the delta-kick perturbation and, after
Fourier transformation, yields the optical conductivity.

The \emph{diamagnetic} current $\Jdia = \Avec(t)|\psi|^2$ is
proportional to the applied vector potential.
In the perturbative kick scheme used in \texttt{RMG},%
\cite{yabana-bertsch-1996,jakowski-rmg-tddft-2025}
the vector potential is a delta-function impulse:
$\Avec(t) = \Avec_0\,\delta(t)$, with $\Avec(t) = 0$ for all
$t > 0$.
Consequently:
\begin{equation}
  \Jdia(\mathbf{r},t>0) = \Avec(t>0)\,|\psi|^2 = 0.
  \label{eq:Jdia-zero}
\end{equation}
The diamagnetic current vanishes identically during the simulation
after the kick.
This is not an approximation: within linear response theory the
complete first-order optical response is encoded in the oscillations
of $\Jpara(t>0)$ and $\mathbf{J}_\mathrm{nl}(t>0)$,%
\cite{yabana-bertsch-1996}
so the macroscopic current that is recorded and Fourier-transformed
is
\begin{equation}
  J_\alpha(t>0)
  = \int_\Omega\!
    \left[\Jpara(\mathbf{r},t)
    + \mathbf{J}_{\mathrm{nl}}(\mathbf{r},t)\right]_\alpha
    \mathrm{d}^3r,
  \label{eq:J-sim}
\end{equation}
where the integral is over the unit cell $\Omega$.

%--------------------------------------------------------------------
\subsubsection{Current Operator for Bloch States:
  the $\mathbf{k}$-Vector Correction}
\label{sec:current-bloch}
%--------------------------------------------------------------------

For a periodic crystal, the Kohn-Sham orbitals satisfy Bloch's
theorem.
We derive how the current operator is modified when expressed in
terms of the cell-periodic part $u_{n,\mathbf{k}}$.

A Bloch state at band index $n$ and crystal momentum $\mathbf{k}$
is written as
\begin{equation}
  \psi_{n,\mathbf{k}}(\mathbf{r},t)
  = e^{\imath\mathbf{k}\cdot\mathbf{r}}\,
    u_{n,\mathbf{k}}(\mathbf{r},t),
  \label{eq:bloch-def}
\end{equation}
where $u_{n,\mathbf{k}}(\mathbf{r}+\mathbf{R},t)
= u_{n,\mathbf{k}}(\mathbf{r},t)$ for all lattice vectors
$\mathbf{R}$.
Applying the canonical momentum operator $-\imath\nabla_\alpha$ to
the full Bloch state and using the product rule:
\begin{align}
  (-\imath\nabla_\alpha)\,\psi_{n,\mathbf{k}}
  &= (-\imath\nabla_\alpha)\!\left[
     e^{\imath\mathbf{k}\cdot\mathbf{r}}\,u_{n,\mathbf{k}}
     \right]
  \nonumber \\
  &= (-\imath)\!\left[
     \imath k_\alpha\,
     e^{\imath\mathbf{k}\cdot\mathbf{r}}\,u_{n,\mathbf{k}}
     + e^{\imath\mathbf{k}\cdot\mathbf{r}}\,
       \nabla_\alpha u_{n,\mathbf{k}}
     \right]
  \nonumber \\
  &= e^{\imath\mathbf{k}\cdot\mathbf{r}}
     \!\left[k_\alpha + (-\imath\nabla_\alpha)\right]
     u_{n,\mathbf{k}}.
  \label{eq:nabla-bloch}
\end{align}
The Bloch phase factor $e^{\imath\mathbf{k}\cdot\mathbf{r}}$
cancels between bra and ket in a matrix element over the unit cell,
so the matrix element of the canonical momentum between two Bloch
states at the same $\mathbf{k}$-point is
\begin{align}
  p^\alpha_{ab}(\mathbf{k})
  &\equiv \langle\psi_{a,\mathbf{k}}|(-\imath\nabla_\alpha)|
    \psi_{b,\mathbf{k}}\rangle
  \nonumber \\
  &= \int_\Omega\!
     e^{-\imath\mathbf{k}\cdot\mathbf{r}}\,u_{a,\mathbf{k}}^*\cdot
     e^{\imath\mathbf{k}\cdot\mathbf{r}}
     \!\left(-\imath\nabla_\alpha + k_\alpha\right)
     u_{b,\mathbf{k}}\,\mathrm{d}^3r
  \nonumber \\
  &= \langle u_{a,\mathbf{k}}|
     \!\left(-\imath\nabla_\alpha + k_\alpha\right)
     |u_{b,\mathbf{k}}\rangle_\Omega.
  \label{eq:para-bloch-me}
\end{align}
The $k_\alpha$ shift arises purely from the Bloch factor and is
absent for $\Gamma$-point ($\mathbf{k} = \mathbf{0}$) calculations.
For finite $\mathbf{k}$-point sampling it must be included to obtain
the correct optical response; its omission would give wrong
intensities proportional to $|\mathbf{k}|^2$ at each $\mathbf{k}$-point.

The full current density summed over all occupied Bloch states is
\begin{equation}
  \mathbf{J}(\mathbf{r},t)
  = \frac{1}{N_k}\!\sum_{n,\mathbf{k}}^{\mathrm{occ}}
    f_{n,\mathbf{k}}\,
    \operatorname{Re}\!\left[
      u_{n,\mathbf{k}}^*\!\left(
      -\imath\nabla + \mathbf{k}
      + \frac{1}{\imath}[\hat{\mathbf{r}},\Vnl]
      \right)u_{n,\mathbf{k}}
    \right],
  \label{eq:J-bloch-total}
\end{equation}
where $N_k$ is the number of $\mathbf{k}$-points and $f_{n,\mathbf{k}}$
are the occupation numbers.
Equation~(\ref{eq:J-bloch-total}) reproduces eq.~(8) of Yabana and
Bertsch\cite{yabana-bertsch-1996} and eq.~(26) of Mattiat and
Luber,\cite{Mattiat-Luber-2022} specialized to the Coulomb gauge
and the velocity gauge with a spatially uniform $\Avec(t)$.

%--------------------------------------------------------------------
\subsubsection{Current Operator Matrix and Density Matrix Formulation}
\label{sec:current-dm-sec}
%--------------------------------------------------------------------

In the density matrix formulation the macroscopic current is
evaluated by contracting the propagated density matrix
$\mathbf{P}_\mathbf{k}(t)$ (Section~\ref{sec:propagation}) with
a fixed current operator matrix that is assembled once before the
time loop.
The current operator matrix element in the Kohn-Sham orbital basis
at $\mathbf{k}$-point $\mathbf{k}$ is
\begin{equation}
  \boxed{
  \Jmat{\alpha}_{ab}(\mathbf{k})
  = \underbrace{
      \langle\psi_{a,\mathbf{k}}|
      (-\imath\nabla_\alpha + k_\alpha)
      |\psi_{b,\mathbf{k}}\rangle
    }_{\displaystyle\texttt{VecPHmatrix()}}
  +
  \underbrace{
      \langle\psi_{a,\mathbf{k}}|
      \tfrac{1}{\imath}[\hat{r}_\alpha,\Vnl]
      |\psi_{b,\mathbf{k}}\rangle
    }_{\displaystyle\texttt{CurrentNlpp()}},
  }
  \label{eq:Jmat-full}
\end{equation}
where the code labels refer to \texttt{RMG} routines described in
Section~\ref{sec:implementation}.
The macroscopic $\alpha$-component of the current is then
\begin{equation}
  J_\alpha(t)
  = \sum_{\mathbf{k}} w_\mathbf{k}\,
    \operatorname{Re}\operatorname{Tr}\!\left[
      \mathbf{P}_\mathbf{k}(t)\cdot
      \boldsymbol{\mathcal{J}}^\alpha_\mathbf{k}
    \right],
  \label{eq:J-dm}
\end{equation}
where $w_\mathbf{k}$ are the $\mathbf{k}$-point weights
(with $\sum_\mathbf{k} w_\mathbf{k} = 1$) and the trace runs over
the active orbital subspace.
In the code, eq.~(\ref{eq:J-dm}) is evaluated at each time step by
a single BLAS complex dot product:
\begin{verbatim}
J[alpha] += Re(zdotc(Norb^2, Pn0, 1, Palpha_matrix, 1)) * kweight;
\end{verbatim}
where \texttt{Palpha\_matrix} contains the stored matrix
$\Jmat{\alpha}_{ab}$ (built by \texttt{VecPHmatrix} +
\texttt{CurrentNlpp} before the time loop) and \texttt{Pn0}
contains the current density matrix elements $P_{ab}(t)$.

%--------------------------------------------------------------------
\subsubsection{Hypervirial Relation and Gauge Invariance}
\label{sec:hypervirial}
%--------------------------------------------------------------------

The hypervirial relation is a consistency condition that connects the
velocity (current) operator in the velocity gauge to the dipole
matrix elements in the length gauge.
It is the periodic-system analog of the relation between oscillator
strengths in the length and velocity forms.

In the Heisenberg picture, the Ehrenfest theorem applied to
$\hat{r}_\alpha$ gives:
\begin{equation}
  \langle\psi_a|\hat{v}_\alpha|\psi_b\rangle
  = \langle\psi_a|\tfrac{1}{\imath}[\hat{r}_\alpha,\hat{H}]
    |\psi_b\rangle
  = (\varepsilon_b - \varepsilon_a)
    \langle\psi_a|\hat{r}_\alpha|\psi_b\rangle,
  \label{eq:ehrenfest}
\end{equation}
where $\varepsilon_a$, $\varepsilon_b$ are the ground-state
Kohn-Sham eigenvalues.
For a Hamiltonian with only local potentials, where $\hat{v}_\alpha
= -\imath\nabla_\alpha$, this gives the standard hypervirial relation
\begin{equation}
  \langle\psi_a|(-\imath\nabla_\alpha)|\psi_b\rangle
  = (\varepsilon_b - \varepsilon_a)
    \langle\psi_a|\hat{r}_\alpha|\psi_b\rangle,
  \label{eq:hypervirial-local}
\end{equation}
which states that the velocity gauge matrix element (left side)
equals the length gauge matrix element (right side): both gauges
are equivalent.

When the Hamiltonian contains a nonlocal pseudopotential, the
velocity operator gains the correction
$\frac{1}{\imath}[\hat{r}_\alpha,\Vnl]$
(eq.~\ref{eq:v-full}).
Inserting into eq.~(\ref{eq:ehrenfest}), and rearranging:
\begin{equation}
  \boxed{
  \langle\psi_a|(-\imath\nabla_\alpha)|\psi_b\rangle
  = (\varepsilon_b - \varepsilon_a)
    \langle\psi_a|\hat{r}_\alpha|\psi_b\rangle
  - \langle\psi_a|\tfrac{1}{\imath}
    [\hat{r}_\alpha,\Vnl]|\psi_b\rangle.
  }
  \label{eq:hypervirial-nlpp}
\end{equation}
Equation~(\ref{eq:hypervirial-nlpp}) is the nonlocal-PP-corrected
hypervirial relation, first given by Starace\cite{starace-1971} and
discussed for RT-TDDFT by Mattiat and
Luber.\cite{Mattiat-Luber-2022}
Its physical meaning is that omitting the commutator term from the
velocity operator would give different results in the velocity and
length gauges — that is, wrong spectra.
Including the correction, and only then, restores the gauge
equivalence and ensures that the computed optical conductivity is
independent of the chosen gauge.
Equation~(\ref{eq:hypervirial-nlpp}) also provides a practical
numerical diagnostic: the left-hand side is computed by
\texttt{VecPHmatrix()}, while the right-hand side can be computed
independently from the Kohn-Sham eigenvalues and the dipole matrix;
agreement between the two, up to the grid discretization error,
validates the implementation.

%--------------------------------------------------------------------
\subsubsection{Sign Convention and Factor of $\imath$ in the Code}
\label{sec:imath-convention}
%--------------------------------------------------------------------

We track the factors of $\imath$ explicitly to verify that the code
implements eq.~(\ref{eq:J-dm}) with the correct sign.

The physical current matrix element
$\Jmat{\alpha}_{ab} = \langle\psi_a|(-\imath\nabla_\alpha + k_\alpha
+ \frac{1}{\imath}[\hat{r}_\alpha,\Vnl])|\psi_b\rangle$
is complex-valued in general.
In \texttt{VecPHmatrix()}, the gradient matrix is computed and stored
multiplied by the prefactor \texttt{I\_t} $= \imath$:
\begin{equation}
  \texttt{Pxmatrix}[ab]
  \leftarrow
  \imath\cdot\langle\psi_a|(-\imath\nabla_\alpha)|\psi_b\rangle
  + \imath\cdot k_\alpha\langle\psi_a|\psi_b\rangle
  = \langle\psi_a|(\nabla_\alpha - \imath k_\alpha)|\psi_b\rangle.
  \label{eq:stored-para}
\end{equation}
That is, \texttt{Pxmatrix} stores $\imath\Jmat{\alpha}_{ab}$ from
the paramagnetic part.
\texttt{CurrentNlpp()} adds the nonlocal contribution with the same
prefactor:
\begin{equation}
  \texttt{Pxmatrix}[ab]
  \mathrel{+}=
  \imath\cdot\langle\psi_a|
  \tfrac{1}{\imath}[\hat{r}_\alpha,\Vnl]|\psi_b\rangle
  = \langle\psi_a|[\hat{r}_\alpha,\Vnl]|\psi_b\rangle.
  \label{eq:stored-nlpp}
\end{equation}
Thus \texttt{Pxmatrix} $= \imath\cdot\Jmat{\alpha}_{ab}$ in total.
The macroscopic current is extracted as
\begin{align}
  J_\alpha
  &= \operatorname{Re}\operatorname{Tr}
     \!\left[\mathbf{P}\cdot\texttt{Pxmatrix}\right]
   = \operatorname{Re}\operatorname{Tr}
     \!\left[\mathbf{P}\cdot(\imath\boldsymbol{\mathcal{J}}^\alpha)\right].
  \label{eq:J-from-code}
\end{align}
For a Hermitian density matrix $\mathbf{P} = \mathbf{P}^\dag$
and an anti-Hermitian current operator matrix
$(\boldsymbol{\mathcal{J}}^\alpha)^\dag = -\boldsymbol{\mathcal{J}}^\alpha$
(which follows from the anti-Hermitian nature of $-\imath\nabla$
and of $\frac{1}{\imath}[\hat{r},\Vnl]$%
\footnote{Both $-\imath\nabla$ and $\frac{1}{\imath}[\hat{r},\Vnl]$
are anti-Hermitian operators when the projectors $\beta^{lm}_I$
are real-valued, which is the case for norm-conserving
pseudopotentials in \texttt{RMG}.}),
the trace satisfies
$\operatorname{Tr}[\mathbf{P}\boldsymbol{\mathcal{J}}^\alpha]
= -\overline{\operatorname{Tr}[\mathbf{P}\boldsymbol{\mathcal{J}}^\alpha]}$,
i.e., it is purely imaginary.
Therefore
\begin{equation}
  J_\alpha
  = \operatorname{Re}\!\left[\imath\operatorname{Tr}
    (\mathbf{P}\boldsymbol{\mathcal{J}}^\alpha)\right]
  = -\operatorname{Im}\operatorname{Tr}
    \!\left[\mathbf{P}\boldsymbol{\mathcal{J}}^\alpha\right]
  = \operatorname{Re}\operatorname{Tr}
    \!\left[\mathbf{P}\cdot\imath\boldsymbol{\mathcal{J}}^\alpha\right],
  \label{eq:J-sign-check}
\end{equation}
which is consistent with eq.~(\ref{eq:J-dm}) and confirms that the
sign convention is correct.
The uniform prefactor \texttt{I\_t} applied in both
\texttt{VecPHmatrix()} and \texttt{CurrentNlpp()} ensures that the
paramagnetic and nonlocal contributions enter the current with
the same relative sign, as required by eq.~(\ref{eq:Jmat-full}).



